% ============ 第3章 §3.5 关于作用函数(原书 pp.265-305)+ 结语(原书 p.305) ============
\section{关于作用函数}\label{sec:3.5}
\index{作用函数}
\index{乘子空间!乘子空间 Mp(Z) 的作用函数}
\index{作用!(乘子空间的情形)}
\index{Rudin--Shapiro 多项式}
\index{Riesz--Thorin 定理}
\index{插值定理!Stein 的}
\index{Hörmander 问题}
\index{Zafran 定理}


在 3.1 节中已给出交换 Banach 代数作用函数的定义,但在那种情形下,是把它作为极大理想空间上的函数空间上的作用来考虑的.本节中的作用函数,例如对交换 Banach 代数 $M(\mathbf{T})$,是通过把它看作真含于极大理想空间的 $\mathbf{Z}$ 上的函数空间来考虑的.这样,本节所处理的交换 Banach 代数的作用函数,均作为对 Fourier 变换上的函数空间的作用来考虑.通过研究函数空间的作用函数,可以看出交换 Banach 代数的特征.

\subsection{作用函数的定义}\label{sec:3.5.1}

\begin{definition}\label{def:3.5.1}
设 $X$ 是局部紧 Hausdorff 空间,$A$,$B$ 是 $X$ 上的函数空间.设 $G$ 为 $\mathbf{C}$ 的子集,$\Phi$ 为 $G$ 上的复值连续函数.称 $\Phi$ 从 $A$ 作用于 $B$,是指:若 $f \in A$,$f(X) \subset G$,则 $\Phi \circ f \in B$.此时称 $\Phi$ 为从 $A$ 到 $B$ 的作用函数.当 $A = B$ 时,简称为 $\Phi$ 作用于 $A$,或 $\Phi$ 是 $A$ 的作用函数.
\end{definition}

一般地,设 $\mathfrak{A}$ 为半单交换 Banach 代数,如 3.1 节所述,通过 Gelfand 表示,$\mathfrak{A}$ 可表示为某个局部紧 Hausdorff 空间 $X$($\mathfrak{A}$ 的极大理想空间)上的函数空间 $A = \widehat{\mathfrak{A}}$.因此,作用函数的研究也可在半单交换 Banach 代数上进行.本节所处理的交换 Banach 代数,是 Fourier 变换恰为 Gelfand 变换的情形.于是,基于 Fourier 变换的性质考虑作为 Fourier 变换的函数空间,并处理其上的作用函数.在这种情况下,借用 $A = \widehat{\mathfrak{A}}$ 的作用函数的说法,也称为 $\mathfrak{A}$ 的作用函数.

\subsection{\texorpdfstring{$A(\mathbf{T})$ 的作用函数}{A(T) 的作用函数}}\label{sec:3.5.2}

下面是 3.1 节定理~\ref{thm:3.1.48} 的特殊情形.但这里利用 $A(\mathbf{T})$ 的性质直接证明.

\begin{theorem}[Wiener--Lévy]\label{thm:3.5.2}
若 $f \in A(\mathbf{T})$,$\Phi$ 在 $f(\mathbf{T})$ 的邻域上解析,则 $\Phi(f) \in A(\mathbf{T})$.
\end{theorem}

\begin{proof}
设 $f$ 为 $\mathbf{T}$ 上的函数.若 $f$ 在 $\mathbf{T}$ 的任意点 $x_{0} \in T$ 的邻域上与 $A(\mathbf{T})$ 的元一致,则 $f \in A(\mathbf{T})$.因为由 $C^{1}(\mathbf{T}) \subset A(\mathbf{T})$,利用 3.1 节系~\ref{cor:3.1.39} 的想法即可容易地证明.其次,设 $\Phi$ 在 $\mathbf{C}$ 的开集 $G$ 上解析.此时,若 $f(\mathbf{T}) \subset G$,$f \in A(\mathbf{T})$,则 $\Phi(f) \in A(\mathbf{T})$.事实上,由前述事实,对 $a \in T = [-\pi, \pi)$,只要在 $a$ 的邻域上存在满足 $F_{a}(x) = \Phi(f(x))$ 的 $F_{a} \in A(\mathbf{T})$ 即可.设 $f(a) = c$.由 $\Phi$ 的解析性,在 $c$ 的邻域上,对适当的 $\varepsilon_{0} > 0$ 可作 Taylor 展开 $\Phi(z) = \sum_{n=0}^{\infty}a_{n}(z - c)^{n}$($|z - c| < \varepsilon_{0}$).在 3.3 节的记号下,把 $\varepsilon > 0$ 取得充分小,则在 $a$ 的适当邻域 $U(a)$ 上,$V_{\varepsilon}^{a}(x)(f(x) - c) = f(x) - c$,$\|V_{\varepsilon}^{a}(x)(f(x) - c)\|_{A(\mathbf{T})} < \varepsilon_{0}$.综上,$\Phi(f(x)) = \sum_{n=0}^{\infty}a_{n}(f(x) - c)^{n}$($x \in U(a)$),$V_{\varepsilon}^{a}(x)(f(x) - c) = f(x) - c$($x \in U(a)$).这里,令 $F_{a}(x) = \sum_{n=0}^{\infty}a_{n}(V_{\varepsilon}^{a}(x)(f(x) - c))^{n}$,则
\begin{equation*}
\sum_{n=0}^{\infty}|a_{n}|\|V_{\varepsilon}^{a}(x)(f(x) - c)\|_{A(\mathbf{T})}^{n} \leq \sum_{n=0}^{\infty}|a_{n}|\varepsilon_{0}^{n} < \infty,
\end{equation*}
故 $\Phi(f(x)) = F_{a}(x)$($x \in U(a)$),$F_{a} \in A(\mathbf{T})$.从而定理得证.
\end{proof}

1958 年,Katznelson \cite{ch3:31} 考虑了上述定理的逆命题,得到如下结果(cf. \cite{ch3:7}, \cite{ch3:8}, \cite{ch3:11}).今天,它被称为 Katznelson 定理.

\begin{theorem}[Katznelson]\label{thm:3.5.3}
设 $\Phi$ 为 $[-1, 1]$ 上的复值函数.假设对任意 $f \in A(\mathbf{T})$,$f(\mathbf{T}) \subset [-1, 1]$,$\Phi(f) \in A(\mathbf{T})$ 成立.此时,$\Phi$ 可延拓为 $[-1, 1]$ 的邻域上的解析函数.
\end{theorem}

证明分若干引理进行.在叙述引理之前,先注意作用函数 $\Phi$ 在 $[-1, 1]$ 上连续.事实上,只需证明在 $x = 0$ 处连续即可,设 $\Phi(0) = 0$ 并假设在 0 处不连续.此时,存在 $\eta, \delta > 0$ 和 $\{a_{n}\}_{n} \subset [-1, 1]$ 使得 $\sum_{n}|a_{n}| < \eta$,$|\Phi(a_{n})| > \delta$.这里,利用 3.3 节定义的函数 $\Delta_{\varepsilon}$,可知存在 $f \in A(\mathbf{T})$,$f(\frac{1}{n}) = a_{n}$($n = 2, 3, \ldots$),$f(\mathbf{T}) \subset [-1, 1]$.此时,因 $\Phi(f) \in A(\mathbf{T})$,$\lim_{n \to \infty}\Phi(f(\frac{1}{n})) = \Phi(f(0))$,可导出与 $\delta$,$\{a_{n}\}_{n}$ 的取法相矛盾.

\begin{lemma}\label{lem:3.5.4}
当 $[-1, 1]$ 上的连续函数 $\Phi$ 作用于 $A(\mathbf{T})$ 时,存在 $\varepsilon, C > 0$,使得实值函数 $f \in A(\mathbf{T})$ 若满足 $\|f\|_{A(\mathbf{T})} < \varepsilon$,则 $\|\Phi(f)\|_{A(\mathbf{T})} \leq C$.
\end{lemma}

\begin{proof}
\begin{sloppypar}
用反证法证明.若命题不成立,则存在实值函数的序列 $\{f_{n}\} \subset A(\mathbf{T})$ 使得 $\|f_{n}\|_{A(\mathbf{T})} \to 0$,$\|\Phi(f_{n})\|_{A(\mathbf{T})} \to \infty$($n \to \infty$).此时,若对任意点 $x \in \mathbf{T}$,存在含 $x$ 的某个闭区间 $I_{x}$ 使得 $\|\Phi(f_{n})\|_{A(I_{x})} \leq C$($C$ 是与 $x, f_{n}$ 无关的常数)(cf. 定义~\ref{def:3.4.9}),则由定理~\ref{thm:3.1.38} 与 $\|\Phi(f_{n})\|_{A(\mathbf{T})} \to \infty$($n \to \infty$)矛盾,故不妨设存在某点 $x_{0} \in \mathbf{T}$,使得对含 $x_{0}$ 的任意闭区间 $I$,$\|\Phi(f_{n})\|_{A(I)} \to \infty$($n \to \infty$).进而,不妨设 $x_{0} = 0$.又由定理~\ref{thm:3.3.1},把闭区间 $I$ 的长度记为 $|I|$,若 $|I| \to 0$,则 $\|\Phi(f_{n})\|_{A(I)} \to |\Phi(f_{n})(0)|$,所以固定 $n$ 时,$\{\|\Phi(f_{n})\|_{A(I)}\}_{I}$ 有界.从而,再次利用定理~\ref{thm:3.1.38} 的想法,对 $\|\Phi(f_{n})\|_{A(I)}$,把含 0 的闭区间 $J \subset I$ 取得足够小,令 $L$ 为 $I \setminus J$ 的闭包,则可做到
\end{sloppypar}
\begin{equation*}
\|\Phi(f_{n})\|_{A(L)} \geq \frac{1}{5}\|\Phi(f_{n})\|_{A(I)} - \|\Phi\|_{\infty}.
\end{equation*}
由此,取适当的数列 $\{n_{k}\}$,$\{n_{k}'\}$,令 $I_{k} = [-2^{-n_{k}}, 2^{-n_{k}}]$,$J_{k} = [-2^{-n_{k}'}, 2^{-n_{k}'}]$,$L_{k} = \overline{I_{k} \setminus J_{k}}$,则可做到 $\|\Phi(f_{n_{k}})\|_{A(L_{k})} \geq k$,$\|f_{n_{k}}\|_{A(\mathbf{T})} \leq 2^{-k}$.进而,令 $V_{k} = 2\Delta_{2^{-n_{k}+1}} - \Delta_{2^{-n_{k}}}$,在 $L_{k}$ 上 $V_{k} - V_{k+1} = 1$;若 $j \neq k$,则可取 $V_{j} - V_{j+1} = 0$.于是令 $f = \sum_{k=5}^{\infty}(V_{k} - V_{k+1})f_{n_{k}}$,由 $\|V_{k}\|_{A(\mathbf{T})} \leq 3$,$\|f\|_{A(\mathbf{T})} \leq 1$.又由 $f(x) = f_{n_{k}}(x)$($x \in L_{k}$),$\|\Phi(f)\|_{A(\mathbf{T})} \geq \|\Phi(f_{n_{k}})\|_{A(L_{k})} \geq k$($k = 5, 6, \ldots$),与 $\Phi(f) \in A(\mathbf{T})$ 矛盾.
\end{proof}

\begin{lemma}\label{lem:3.5.5}
若 $g \in A(\mathbf{T})$,$h \in A(\mathbf{T})$,则
\begin{equation*}
\lim_{k \to \infty}\|g(x)h(kx)\|_{A(\mathbf{T})} = \|g\|_{A(\mathbf{T})}\|h\|_{A(\mathbf{T})}
\end{equation*}
成立.
\end{lemma}

\begin{proof}
\begin{sloppypar}
$\lim_{k \to \infty}\|g(x)h(kx)\|_{A(\mathbf{T})} \leq \|g\|_{A(\mathbf{T})}\|h\|_{A(\mathbf{T})}$ 是显然的,故证明反向不等式.设 $h_{n}(x) = \sum_{|j| \leq n}\widehat{h}(j)e^{ijx}$,则对任意 $\varepsilon > 0$,存在 $N$ 使得 $\|h\|_{A(\mathbf{T})} - \varepsilon < \|h_{N}\|_{A(\mathbf{T})}$.此时,由 $g(x)h(kx) = g(x)h_{N}(kx) + g(x)(h(kx) - h_{N}(kx))$,
$\|g(x)h(kx)\|_{A(\mathbf{T})} \geq \|g(x)h_{N}(kx)\|_{A(\mathbf{T})} - \|g(x)(h(kx) - h_{N}(kx))\|_{A(\mathbf{T})}$,且
\end{sloppypar}
\begin{align*}
\|g(x)(h(kx) - h_{N}(kx))\|_{A(\mathbf{T})} &\leq \|g\|_{A(\mathbf{T})}\|h(kx) - h_{N}(kx)\|_{A(\mathbf{T})}\\
&\leq \varepsilon\|g\|_{A(\mathbf{T})}
\end{align*}
成立.因此
\begin{equation*}
\|g(x)h(kx)\|_{A(\mathbf{T})} \geq \|g(x)h_{N}(kx)\|_{A(\mathbf{T})} - \varepsilon\|g\|_{A(\mathbf{T})}
\end{equation*}
成立.又与 $\{h_{n}\}_{n}$ 同样地定义 $\{g_{n}\}_{n}$,则因 $g(x)h_{N}(kx) = g_{M}(x)h_{N}(kx) + (g(x) - g_{M}(x))h_{N}(kx)$,把 $M$ 取得足够大时,
\begin{equation*}
\|g(x)h_{N}(kx)\|_{A(\mathbf{T})} \geq \|g_{M}(x)h_{N}(kx)\|_{A(\mathbf{T})} - \varepsilon\|h\|_{A(\mathbf{T})}
\end{equation*}
成立.这里把 $k$ 取得足够大时,可证明
\begin{equation*}
\|g_{M}(x)h_{N}(kx)\|_{A(\mathbf{T})} \geq (\|g\|_{A(\mathbf{T})} - 2\varepsilon)(\|h\|_{A(\mathbf{T})} - 2\varepsilon),
\end{equation*}
故对 $\eta > 0$,只要 $k$ 足够大,
\begin{equation*}
\|g(x)h(kx)\|_{A(\mathbf{T})} \geq \|g\|_{A(\mathbf{T})}\|h\|_{A(\mathbf{T})} - \eta,
\end{equation*}
命题得证.
\end{proof}

\begin{lemma}\label{lem:3.5.6}
对 $\varepsilon > 0$,
\begin{equation*}
\|e^{i\varepsilon\cos x}\|_{A(\mathbf{T})} = 1 + \varepsilon + O(\varepsilon^{2})
\end{equation*}
成立.这里,$g(x) = O(x^{2})$($x \to 0$)表示 $g(x)/x^{2}$ 在 $x \to 0$ 时有界.
\end{lemma}

\begin{proof}
在
\begin{equation*}
e^{i\varepsilon\cos x} = \sum_{n=0}^{\infty}\frac{(i\varepsilon\cos x)^{n}}{n!}
\end{equation*}
中,由 $\|\cos x\|_{A(\mathbf{T})} \leq 1$ 得 $\|e^{i\varepsilon\cos x}\|_{A(\mathbf{T})} \leq 1 + \varepsilon + O(\varepsilon^{2})$.反之,由 $\|e^{i\varepsilon\cos x}\|_{A(\mathbf{T})} \geq 1 + \varepsilon - \left\|\sum_{n=2}^{\infty}\frac{(i\varepsilon\cos x)^{n}}{n!}\right\|_{A(\mathbf{T})}$ 及 $\|\cos x\|_{A(\mathbf{T})} \leq 1$ 得 $\|e^{i\varepsilon\cos x}\|_{A(\mathbf{T})} \geq 1 + \varepsilon - O(\varepsilon^{2})$.
\end{proof}

\begin{lemma}\label{lem:3.5.7}
对 $r > 0$,
\begin{equation*}
\sup\left\{\|e^{if}\|_{A(\mathbf{T})}: f \text{ 是实值函数, } \|f\|_{A(\mathbf{T})} \leq r\right\} = e^{r}
\end{equation*}
成立.
\end{lemma}

\begin{proof}
若 $\|f\|_{A(\mathbf{T})} \leq r$,则 $\|e^{if}\|_{A(\mathbf{T})} \leq e^{r}$,故证明反向不等式.令 $k_{1} = 1$,设 $f(x) = \frac{r}{n}(\cos k_{1}x + \cdots + \cos k_{n}x)$($k_{2}, \ldots, k_{n}$ 是自然数),由引理~\ref{lem:3.5.5},把 $k_{2}, \ldots, k_{n}$ 取得足够大,则
\begin{equation*}
\left\|\prod_{j=1}^{n}e^{i\frac{r}{n}\cos k_{j}x}\right\|_{A(\mathbf{T})} > \left(1 - \frac{1}{n}\right)\prod_{j=1}^{n}\left\|e^{i\frac{r}{n}\cos k_{j}x}\right\|_{A(\mathbf{T})}.
\end{equation*}
进而利用引理~\ref{lem:3.5.6},对任意 $\varepsilon > 0$,把 $n$ 取得足够大,则 $\left(1 - \frac{1}{n}\right)\left(1 + \frac{r}{n} + O\left(\frac{r^{2}}{n^{2}}\right)\right)^{n} > (1 - \varepsilon)e^{r}$.此时,得 $\|f\|_{A(\mathbf{T})} \leq r$ 且 $\|e^{if}\|_{A(\mathbf{T})} > (1 - \varepsilon)e^{r}$.
\end{proof}

\begin{proof}[定理~\ref{thm:3.5.3} 的证明]
如前所述,先注意 $\Phi$ 在 $[-1, 1]$ 上连续.首先证明 $\Phi(x)$ 在 $x = 0$ 的邻域上解析.由引理~\ref{lem:3.5.4},只要 $\alpha > 0$ 足够小,就存在常数 $C > 0$,使得当 $\|f\|_{A(\mathbf{T})} \leq 1$,$f$ 为实值函数,$-\pi \leq x < \pi$ 时,$\|\Phi(\alpha\sin(f(t) + x))\|_{A(\mathbf{T})} \leq C$.现在,令 $\Phi_{1}(x) = \Phi(\alpha\sin x)$,则由向量值函数的积分,对实值函数 $f \in A(\mathbf{T})$,
\begin{equation*}
\widehat{\Phi_{1}}(n)e^{inf(t)} = \frac{1}{2\pi}\int_{-\pi}^{\pi}\Phi_{1}(f(t) + x)e^{-inx}\dd x
\end{equation*}
成立.从而得 $\left|\widehat{\Phi_{1}}(n)\right|\|e^{inf}\|_{A(\mathbf{T})} \leq C$.此时,由引理~\ref{lem:3.5.7},$\left|\widehat{\Phi_{1}}(n)\right| \leq Ce^{-|n|}$,所以由众所周知的事实,$\Phi_{1}(z)$ 在 $|\mathrm{Im}z| < 1$ 上解析,利用 $\Phi(x) = \Phi_{1}\left(\arcsin\frac{x}{\alpha}\right)$,$\Phi$ 在 $x = 0$ 的邻域上解析.若 $a \in (-1, 1)$,令 $\Psi(x) = \Phi(x + a)$,取适当的 $\varepsilon > 0$,则 $\Psi$ 是在 $[-\varepsilon, \varepsilon]$ 上定义的 $A(\mathbf{T})$ 的作用函数.因此在 $x = 0$ 的邻域上解析,所以 $\Phi$ 在 $x = a$ 的邻域上解析.最后证明在 $x = \pm 1$(端点)处解析.令 $\Psi(x) = \Phi(1 - x^{2})$($x \in [-1, 1]$),因 $\Psi$ 是在 $[-1, 1]$ 上定义的 $A(\mathbf{T})$ 的作用函数,在 $x = 0$ 的邻域上解析.因 $\Psi$ 是偶函数,
\begin{equation*}
\Psi(x) = \sum_{n=0}^{\infty}a_{n}x^{2n}.
\end{equation*}
因此,$\Phi(1 - x) = \sum_{n=0}^{\infty}a_{n}x^{n}$,故对适当的 $\varepsilon_{0} > 0$,
\begin{equation*}
\Phi(x) = \sum_{n=0}^{\infty}a_{n}(1 - x)^{n} \quad (|x| < \varepsilon_{0})
\end{equation*}
可作 Taylor 展开.从而在 $x = 1$ 的邻域上解析.同样可证明在 $x = -1$ 的邻域上解析.
\end{proof}

\subsection{\texorpdfstring{$L^1(\mathbf{T})$ 与 $M(\mathbf{T})$ 的作用函数}{L1(T) 与 M(T) 的作用函数}}\label{sec:3.5.3}

这里参照 Rudin \cite{ch3:11}(cf. \cite{ch3:11} 的参考文献)叙述 $L^{1}(\mathbf{T})$ 与 $M(\mathbf{T})$ 的作用函数.

\begin{definition}\label{def:3.5.8}
$[-1, 1]$ 上的函数 $\Phi(x)$ 是 $L^{1}(\mathbf{T})$ 的作用函数,是指:令 $A(\mathbf{Z}) = \{\{\widehat{f}(n)\}_{n \in \mathbf{Z}}: f \in L^{1}(\mathbf{T})\}$,$\Phi$ 是 $A(\mathbf{Z})$ 上的作用函数.即,当 $\widehat{f} \in A(\mathbf{Z})$,$\widehat{f}(\mathbf{Z}) \subset [-1, 1]$ 时,存在 $\Phi \circ f \in L^{1}(\mathbf{T})$ 使得 $\widehat{\Phi \circ f} = \Phi(\widehat{f}) \in A(\mathbf{Z})$.又,$[-1, 1]$ 上的函数 $\Phi(x)$ 是 $M(\mathbf{T})$ 的作用函数,是指:令 $B(\mathbf{Z}) = \{\{\widehat{\mu}(n)\}_{n \in \mathbf{Z}}: \mu \in M(\mathbf{T})\}$,$\Phi$ 是 $B(\mathbf{Z})$ 上的作用函数.即,当 $\widehat{\mu} \in B(\mathbf{Z})$,$\widehat{\mu}(\mathbf{Z}) \subset [-1, 1]$ 时,存在 $\Phi \circ \mu \in M(\mathbf{T})$ 使得 $\widehat{\Phi \circ \mu} = \Phi(\widehat{\mu}) \in B(\mathbf{Z})$.把 $A(\mathbf{Z})$ 看作由 $L^{1}(\mathbf{T})$ 自然引入的范数及作为函数的通常运算所构成的交换 Banach 代数,$B(\mathbf{Z})$ 同样看作由 $M(\mathbf{T})$ 引入的交换 Banach 代数.
\end{definition}

\begin{theorem}\label{thm:3.5.9}
设 $\Phi$ 为在 $[-1, 1]$ 上定义的函数.此时,$\Phi$ 是 $L^{1}(\mathbf{T})$ 的作用函数的充要条件是:$\Phi(0) = 0$ 且 $\Phi(x)$ 在 0 的邻域上解析.
\end{theorem}

该定理的充分条件如下证明.设 $\Phi$ 在 $|x| < \varepsilon$($\varepsilon > 0$)上解析.取 $f \in L^{1}(\mathbf{T})$,$\hat{f}(\mathbf{Z}) \subset [-1, 1]$,则对任意 $\varepsilon > 0$,取适当的三角多项式 $P$,$Q$,可使 $\Phi(\widehat{f}(n)) = \widehat{Q}(n) + \Phi(\widehat{f}(n) - \widehat{P}(n))$($n \in \mathbf{Z}$),$\|f - P\|_{1} < \varepsilon$.因此由 $\Phi$ 的假设,$\Phi(\widehat{f}(n) - \widehat{P}(n)) \in A(\mathbf{Z})$,故 $\Phi$ 是 $L^{1}(\mathbf{T})$ 的作用函数.

必要性的证明如下进行.首先,因常数函数是 $L^{1}(\mathbf{T})$ 的元,若 $\Phi$ 作用于 $L^{1}(\mathbf{T})$,由 Riemann--Lebesgue 引理得 $\Phi(0) = 0$.

其次,叙述两个引理.

\begin{lemma}\label{lem:3.5.10}
$\Phi(x)$ 在 $x = 0$ 处连续.
\end{lemma}

\begin{proof}
设 $0 < \eta < 1$.若 $\Phi$ 在 0 处不连续,则存在 $\{a_{n}\} \subset R$,$\delta > 0$ 使得 $\sum_{n=1}^{\infty}2|a_{n}| < \eta$,$|\Phi(a_{n})| > \delta$.此时,令 $k_{n} = 5^{n}$,$f(x) = \sum_{n=1}^{\infty}2a_{n}\cos k_{n}x$,则 $\Phi(\widehat{f}) \in A(\mathbf{Z}) \subset B(\mathbf{Z})$.因此存在 $\mu \in M(\mathbf{T})$ 使得
\begin{equation*}
\widehat{\mu}(l) = \left\{\begin{array}{ll}
\Phi(a_{n}) & (l = \pm k_{n})\\
0 & (l \neq \pm k_{n})
\end{array}\right.
\end{equation*}
于是,令 $d\mu_{n}(x) = e^{-ik_{n}x}\dd\mu(x)$,设 $\sigma \in M(\mathbf{T})$ 为 $\{\mu_{n}\} \subset M(\mathbf{T})$ 关于 $*$-弱拓扑的一个聚点.由 $|\widehat{\mu_{n}}(0)| = |\Phi(a_{n})| > \delta$,$\sigma \neq 0$.又对 $l \neq 0$,$l \in \mathbf{Z}$,$\widehat{\mu_{n}}(l) = \widehat{\mu}(k_{n} + l)$,故注意 $\widehat{\mu_{n}}(l) \neq 0$ 只在 $l = -2k_{n}$ 时发生.因此 $\widehat{\sigma}(l) = 0$($l \neq 0$,$l \in \mathbf{Z}$).从而 $\sigma$ 是非零的常数函数.可是,由 Riemann--Lebesgue 引理,对任意三角多项式 $f(x)$,
\begin{equation*}
\int_{0}^{2\pi}f(x)\dd\sigma(x) = \lim_{n \to \infty}\int_{0}^{2\pi}f(x)\dd\mu_{n}(x) = \lim_{n \to \infty}\int_{0}^{2\pi}f(x)e^{-ik_{n}x}\dd\nu(x)
\end{equation*}
成立.其中,$\nu$ 是 $\mu$ 集中于 Lebesgue 测度为零的集合上的部分(奇异测度部分).由此,利用 $\sigma, \nu \in M(\mathbf{T})$,对可测集 $E$ 有 $|\sigma|(E) \leq |\nu|(E)$.这与 $\sigma$ 是非零常数函数矛盾.
\end{proof}

\begin{lemma}\label{lem:3.5.11}
设 $\Phi$ 为 $[-\eta, \eta]$($0 < \eta < 1$)上的,作用于 $L^{1}(\mathbf{T})$ 的函数.此时,存在 $0 < \varepsilon_{0} < 1$,$C > 0$,使得若 $\|\sigma\|_{M(\mathbf{T})} < \varepsilon_{0}$,$\widehat{\sigma}(\mathbf{Z}) \subset \mathbf{R}$,则 $\Phi \circ \widehat{\sigma} \in B(\mathbf{Z})$,$\|\Phi \circ \sigma\|_{M(\mathbf{T})} \leq C$.其中设 $\widehat{\Phi \circ \sigma}(n) = \Phi(\widehat{\sigma}(n))$.
\end{lemma}

\begin{proof}
只需对三角多项式 $f$(代替引理中的 $\sigma$)证明即可.事实上,设对三角多项式已证明.设 $\|\sigma\|_{M(\mathbf{T})} < \varepsilon_{0}$,$\widehat{\sigma}(\mathbf{Z}) \subset \mathbf{R}$.此时,取三角多项式列 $\{k_{n}\}$ 满足 $\hat{k_{n}}(m) = 1$($|m| \leq n$),$\|k_{n}\|_{1} = 1$,则 $\|\Phi \circ (\sigma * k_{n})\|_{M(\mathbf{T})} \leq C$.$\{\Phi \circ (\sigma * k_{n})\}$ 关于 $*$-弱拓扑的极限是 $\Phi \circ \sigma$,故得结果.于是,假设存在三角多项式列 $\{f_{n}\}$ 使得 $\widehat{f_{n}}(\mathbf{Z}) \subset [-1, 1]$,$\|f_{n}\|_{1} < 2^{-n}$,$\|\Phi \circ f_{n}\|_{M(\mathbf{T})} > n$($n \in \mathbf{N}$).这里,选取 $\mathbf{T}$ 上的三角多项式 $h_{n}$ 作为三角多项式列 $\{h_{n}\}$:$\|h_{n}\|_{1} < 2$,在 $\operatorname{supp}\widehat{f_{n}}$ 上为 1,且 $0 \leq \widehat{h_{n}} \leq 1$.此时,注意 Fourier 系数的平移不改变函数的范数以及 $\Phi(0) = 0$,不妨设 $\operatorname{supp}\widehat{h_{n}} \cap \operatorname{supp}\widehat{h_{m}} = \emptyset$($n \neq m$),并令 $g = \sum_{n=1}^{\infty}f_{n}$,则 $g \in L^{1}(\mathbf{T})$,且 $h_{n} * (\Phi \circ g) = h_{n} * (\Phi \circ f_{n}) = \Phi \circ f_{n}$,$\|\Phi \circ f_{n}\|_{1} = \|h_{n} * (\Phi \circ g)\|_{1} \leq 2\|\Phi \circ g\|_{M(\mathbf{T})} < \infty$.这与 $\|\Phi \circ f_{n}\|_{M(\mathbf{T})} \to \infty$($n \to \infty$)矛盾.
\end{proof}

\begin{proof}[定理~\ref{thm:3.5.9} 的证明]
先注意 $\Phi(0) = 0$.设 $\Phi$ 作用于 $A(\mathbf{Z})$.对 $\varepsilon_{0} > 0$ 和 $a$($|a| < \varepsilon_{0}$),令 $\Phi_{1}(x) = \Phi(x + a) - \Phi(a)$,由引理~\ref{lem:3.5.11},当 $\widehat{f} \in A(\mathbf{Z})$,$\widehat{f}(\mathbf{Z}) \subset [-\eta, \eta]$($0 < \eta < 1$),$\|f\|_{1} < \varepsilon_{0} - |a|$ 时,$\Phi_{1}(\widehat{f}) \in B(\mathbf{Z})$.此时,利用引理~\ref{lem:3.5.10} 的想法,$\Phi_{1}$ 在 $x = 0$ 处连续.因此 $\Phi$ 在 $x = a$ 处连续,得 $\Phi$ 在 $(-\varepsilon_{0}, \varepsilon_{0})$ 上连续.于是,取 $0 < r < \varepsilon_{0}/e$,令 $\Phi_{2}(x) = \Phi(r\sin x)$($x \in \mathbf{R}$),则 $\Phi_{2}$ 是 $\mathbf{R}$ 上的 $2\pi$ 周期连续函数.现在,取 $\mu \in M(\mathbf{T})$ 满足 $\|\mu\|_{M(\mathbf{T})} \leq 1$,$\widehat{\mu}(\mathbf{Z}) \subset \mathbf{R}$,则对 $a \in \mathbf{R}$,
\begin{align*}
&\|\sin(\widehat{\mu} + a)\|_{B(\mathbf{Z})} \leq |\cos a|\|\sin\widehat{\mu}\|_{B(\mathbf{Z})} + |\sin a|\|\cos\widehat{\mu}\|_{B(\mathbf{Z})}\\
&\qquad \leq e^{\|\widehat{\mu}\|_{B(\mathbf{Z})}} = e^{\|\mu\|_{M(\mathbf{T})}} \leq e
\end{align*}
成立.因此,若 $a \in R$,$\widehat{\mu}(\mathbf{Z}) \subset \mathbf{R}$,$\|\mu\|_{M(\mathbf{T})} \leq 1$,则由 $\|r\sin(\widehat{\mu} + a)\|_{B(\mathbf{Z})} \leq re \leq \varepsilon_{0}$,利用引理~\ref{lem:3.5.11} 得 $\|\Phi_{2}(\widehat{\mu} + a)\|_{B(\mathbf{Z})} \leq \|\Phi(r\sin(\widehat{\mu} + a))\|_{B(\mathbf{Z})} \leq C$.这里,因 $\Phi_{2}$ 是连续函数,
\begin{equation*}
\Phi_{2}(x) \sim \sum_{-\infty}^{\infty}c_{n}e^{inx}
\end{equation*}
作 Fourier 级数展开.取注意~\ref{note:3.3.12} 所述的,作为 Cantor 集的 Kronecker 集 $K$,设 $\mu \in M(K \cup (-K))$ 满足任意 1 点的测度为 0,$\widehat{\mu}(\mathbf{Z}) \subset \mathbf{R}$,$0 < \|\mu\|_{M(\mathbf{T})} \leq 1$.因 $\Phi_{2}$ 连续,利用向量值函数的积分(cf. \cite{ch3:3}, \cite{ch3:8}),估计
\begin{equation*}
c_{n}e^{in\widehat{\mu}(j)} = \lim_{N \to \infty}\frac{1}{N}\sum_{k=1}^{N}\Phi_{2}\left(\frac{2\pi k}{N} + \widehat{\mu}(j)\right)e^{-i2\pi k/N},
\end{equation*}
得 $|c_{n}|\|e^{in\mu}\|_{M(\mathbf{T})} \leq C$($n \in \mathbf{Z}$).由 $\mu$ 的性质,$\|e^{in\mu}\|_{M(\mathbf{T})} = e^{|n|\|\mu\|_{M(\mathbf{T})}}$($n \in \mathbf{Z}$),故得 $|c_{n}| \leq Ce^{-|n|}$($n \in \mathbf{Z}$).从而 $\Phi_{2}(z)$ 在 $|\mathrm{Im}z| < 1$ 上解析.因此 $\Phi(x) = \Phi_{2}\left(\arcsin\frac{x}{r}\right)$ 在 0 的邻域上解析.
\end{proof}

关于 $M(\mathbf{T})$ 的作用函数,下面的事实成立.

\begin{theorem}\label{thm:3.5.12}
$[-1, 1]$ 上的复值函数 $\Phi$ 是 $M(\mathbf{T})$ 的作用函数的充要条件是:$\Phi$ 可延拓为整函数.
\end{theorem}

首先证明下面的事实.

\begin{lemma}\label{lem:3.5.13}
若 $\Phi_{1}$ 是 $\mathbf{R}$ 上的 $2\pi$ 周期连续函数且作用于 $B(\mathbf{Z})$,则当 $\widehat{\mu} \in B(\mathbf{Z})$,$\widehat{\mu}(\mathbf{Z}) \subset \mathbf{R}$ 时,存在 $\delta > 0$,$C < \infty$ 使得
\begin{equation*}
\|\Phi_{1}(\widehat{\mu} + \widehat{\sigma})\|_{B(\mathbf{Z})} \leq C \quad (\sigma \in M(\mathbf{T}),\ \widehat{\sigma} \in \mathbf{R},\ \|\sigma\|_{M(\mathbf{T})} < \delta).
\end{equation*}
\end{lemma}

\begin{proof}
不妨设 $\Phi_{1}(0) = 0$.只需证明 $\|\Phi_{1}(\widehat{\mu} + \widehat{f})\|_{B(\mathbf{Z})} \leq C$($f \in L^{1}(\mathbf{T})$,$\widehat{f} \in \mathbf{R}$,$\|f\|_{1} < \delta$).因为,由 $L^{1}(\mathbf{T})$ 是 $M(\mathbf{T})$ 的理想以及 $M(\mathbf{T})$ 的闭单位球关于 $*$-弱拓扑是紧集即可证明.现在,假设上述事实不成立.此时,存在 $\mathbf{T}$ 上取实值的三角多项式列 $\{f_{n}\}$ 使得 $\|f_{n}\|_{1} \to 0$,$\|\Phi_{1} \circ (\widehat{\mu} + \widehat{f_{n}})\|_{B(\mathbf{Z})} \to \infty$($n \to \infty$).为导出矛盾,如下构造 $\{\lambda_{j}\} \subset M(\mathbf{T})$.令 $n_{1} = 1$,设 $n_{1}, \ldots, n_{j}$ 与三角多项式列 $f_{n_{1}}, \ldots, f_{n_{j}}$ 已选出,令 $\lambda_{j} = \mu + f_{n_{1}} + \cdots + f_{n_{j}}$.此时,取次数足够高的 Fejér 核 $k_{j}(x) = \sum_{|m| \leq n_{j}}\left(1 - \frac{|m|}{n_{j} + 1}\right)e^{imx}$,则
\begin{equation*}
\|k_{j} * (\Phi_{1} \circ \lambda_{j})\|_{M(\mathbf{T})} > \frac{1}{2}\|\Phi_{1} \circ \lambda_{j}\|_{M(\mathbf{T})}.
\end{equation*}
因 $\Phi_{1}$ 连续,把 $n_{j+1}$ 取得足够大,可做到
\begin{align*}
\|f_{n_{j+1}}\|_{1} &< 2^{-j}, \quad \|\Phi_{1} \circ (\lambda_{j} + f_{n_{j+1}})\|_{M(\mathbf{T})} \geq j + 1,\\
\left\|k_{i} * \left((\Phi_{1} \circ \lambda_{j}) - \Phi_{1} \circ (\lambda_{j} + f_{n_{j+1}})\right)\right\|_{1} &< 2^{-j} \quad (1 \leq i \leq j).
\end{align*}
这样归纳地取 $\{n_{j}\}$ 与 $\{k_{j}\}$ 的序列,令 $\lambda = \mu + \sum_{j=1}^{\infty}f_{n_{j}}$,$\tau = \Phi_{1} \circ \lambda$.这里注意 $\widehat{\tau}(\gamma) = \lim_{N \to \infty}\Phi_{1}(\widehat{\lambda_{N}}(\gamma))$,考虑到 $k_{s}$ 是 Fejér 核,
\begin{equation*}
\lim_{N \to \infty}\|k_{s} * (\tau - \Phi_{1} \circ \lambda_{N})\|_{M(\mathbf{T})} = 0 \quad (s = 1, 2, \ldots)
\end{equation*}
成立.因此
\begin{equation*}
k_{s} * (\tau - \Phi_{1} \circ \lambda_{s}) = \lim_{N \to \infty}\sum_{j=s}^{N-1}k_{s} * (\Phi_{1} \circ \lambda_{j+1} - \Phi_{1} \circ \lambda_{j}),
\end{equation*}
所以
\begin{equation*}
\|k_{s} * (\tau - \Phi_{1} \circ \lambda_{s})\|_{M(\mathbf{T})} \leq \sum_{j=s}^{\infty}2^{-j} \leq 1 \quad (s \in \mathbf{N}).
\end{equation*}
从而
\begin{align*}
2\|\tau\|_{M(\mathbf{T})} &\geq \|k_{s} * \tau\|_{M(\mathbf{T})}\\
&\geq \|k_{s} * (\Phi_{1} \circ \lambda_{s})\|_{M(\mathbf{T})} - 1\\
&\geq \frac{1}{2}\|\Phi_{1} \circ \lambda_{s}\|_{M(\mathbf{T})} - 1 \geq \frac{1}{2}s - 1.
\end{align*}
这是矛盾.
\end{proof}

\begin{proof}[定理~\ref{thm:3.5.12} 的证明]
只需证明必要条件.取 $0 < r < 1$,令 $\Phi_{3}(x) = \Phi(r\sin x)$,则 $\Phi_{3}$ 是在 $\mathbf{R}$ 上定义的 $2\pi$ 周期函数,且是 $M(\mathbf{T})$ 的作用函数.此时,利用引理~\ref{lem:3.5.10} 可证明 $\Phi_{3}$ 是 $\mathbf{R}$ 上的连续函数.又,把 $\Phi_{3}$ 的 Fourier 级数记为 $\Phi_{3} \sim \sum_{n}c_{n}e^{inx}$.此时,与定理~\ref{thm:3.5.9} 的证明一样利用引理~\ref{lem:3.5.13},取适当的 $\mu \in M(\mathbf{T})$(cf. 注意~\ref{note:3.3.12}),则
\begin{equation*}
|c_{n}| \leq Ce^{-|n|\|\mu\|_{M(\mathbf{T})}} \quad (n \in \mathbf{Z})
\end{equation*}
成立.由此,$\Phi_{3}(z)$ 在 $|\mathrm{Im}z| < \|\mu\|_{M(\mathbf{T})}$ 上解析,且 $\|\mu\|_{M(\mathbf{T})}$ 可以取得任意大,故 $\Phi_{3}$ 与某个整函数一致.现在,令 $\Psi_{j}(x) = \Phi(r_{j}\sin x)$($j = 1, 2$,$0 < r_{1} < r_{2} < 1$),如前所述,$\Psi_{j}$ 在 $\mathbf{R}$ 上与某个整函数一致.由此,对 $j = 1, 2$,
\begin{equation*}
\Phi(x) = \Psi_{j}\left(\arcsin\frac{x}{r_{j}}\right) \quad (-r_{j} < x < r_{j}),
\end{equation*}
$\Phi(x)$ 在 0 的邻域上与解析函数一致.又,利用多值函数 $\arcsin$ 作解析延拓,可证明 $\Phi$ 在 $[-1, 1]$ 上与某个整函数一致.
\end{proof}

\begin{note}\label{note:3.5.14}
由定理~\ref{thm:3.5.12} 也可导出 $M(\mathbf{T})$ 的 Wiener--Pitt 现象.事实上,令 $\Phi(z) = \frac{1}{1 + z^{2}}$,因该函数不是整函数,故不作用于 $B(\mathbf{Z})$.因此,可知存在 $\mu \in M(\mathbf{T})$ 使得 $\hat{\mu}(n) \in \mathbf{R}$($n \in \mathbf{Z}$),$i \in \mathrm{Sp}(\mu, M(\mathbf{T}))$.
\end{note}

\subsection{\texorpdfstring{乘子空间 $M_p(\mathbb{Z})$ 的作用函数}{乘子空间 M_p(Z) 的作用函数}}\label{sec:3.5.4}

这里叙述乘子空间 $M_{p}(\mathbf{Z})$ 的作用函数.$[-1, 1]$ 上的函数 $\Phi(x)$ 是 $M_{p}(\mathbf{Z})$ 上的作用函数,是指:当 $\hat{T} \in M_{p}(\mathbf{Z})$,$\hat{T}(\mathbf{Z}) \subset [-1, 1]$ 时,存在 $\Phi \circ T \in \mathfrak{M}_{p}(\mathbf{T})$ 使得 $\widehat{\Phi \circ T} = \Phi(\hat{T}) \in M_{p}(\mathbf{Z})$.如 3.2 节所见,$\mathfrak{M}_{p}(\mathbf{T})$ 关于算子范数成为交换 Banach 代数.因此,$M_{p}(\mathbf{Z}) = \{\hat{T} \in \ell^{\infty}: T \in \mathfrak{M}_{p}(\mathbf{T})\}$ 关于由 $\mathfrak{M}_{p}(\mathbf{T})$ 引入的范数与函数的运算成为含单位元的交换 Banach 代数.如前所见,$\mathfrak{M}_{1}(\mathbf{T}) = M(\mathbf{T})$(或 $M_{1}(\mathbf{Z}) = B(\mathbf{Z})$),$M_{2}(\mathbf{Z}) = \ell^{\infty}$,$\mathfrak{M}_{1}(\mathbf{T}) \subset \mathfrak{M}_{p}(\mathbf{T}) \subset \mathfrak{M}_{2}(\mathbf{T})$($1 \leq p \leq 2$)成立.从而,当 $T \in \mathfrak{M}_{p}(\mathbf{T})$ 时,$\overline{\hat{T}(\mathbf{Z})} \subset \mathrm{Sp}(T, \mathfrak{M}_{p}(\mathbf{T}))$.特别是,若 $f \in L^{1}(\mathbf{T})$,则关于算子 $T_{f}(g) = f * g$ 有 $\mathrm{Sp}(T_{f}, \mathfrak{M}_{p}(\mathbf{T})) = \hat{f}(\mathbf{Z}) \cup \{0\}$($1 \leq p \leq 2$).

1960 年,Hörmander \cite{ch3:22} 把 $m_p$ 取为由 $\{T_f : f \in L^1(\mathbf{T})\}$ 在 $\mathfrak{M}_p(\mathbf{T})$ 中的闭包构成的交换 Banach 代数,证明了下面的事实:

\begin{proposition}\label{prop:3.5.15}
(1) $\Delta(m_p) = \mathbf{Z}$.

(2) 在 0 的邻域上解析的函数 $\Phi$ 若满足 $\phi(0) = 0$,则作用于 $m_{p}$.

(3) 对 $1 < p < 2$,令 $C_{0}\mathfrak{M}_{p}(\mathbf{T}) = \{T \in \mathfrak{M}_{p}(\mathbf{T}): \hat{T}(n) \to 0 \ (|n| \to \infty)\}$.设 $1 < p < q < 2$,若满足 $\Phi(0) = 0$ 的函数 $\Phi$ 在 0 的邻域上解析,则 $\Phi$ 从 $C_{0}\mathfrak{M}_{p}(\mathbf{T})$ 作用于 $\mathfrak{M}_{q}(\mathbf{T})$.
\end{proposition}

又,他还提出了如下问题:``当 $1 < p < 2$ 时,$C_{0}\mathfrak{M}_{p}(\mathbf{T}) = m_{p}$ 是否成立?''Igari \cite{ch3:25} 对命题~\ref{prop:3.5.15}(3) 的 $p = q$ 的情形作了相应研究,决定了 $\mathfrak{M}_{p}(\mathbf{T})$($= M_{p}(\mathbf{Z})$)的作用函数.这里叙述 S. Igari 的结果.另外,Hörmander 问题由 M. Zafran 解决一事,将在 3.5.6 小节补充.

下面的插值定理省略证明(cf. \cite{ch3:4}).

\begin{theorem}[Riesz--Thorin]\label{thm:3.5.16}
设 $(X, \mu)$,$(Y, \nu)$ 为 $\sigma$-有限的测度空间.设 $1 \leq p_{0}, q_{0}, p_{1}, q_{1} \leq \infty$($p_{0} \neq p_{1}$,$q_{0} \neq q_{1}$),$T$ 为从 $X$ 上简单函数的空间到 $Y$ 上可测函数空间的线性算子,且满足:存在常数 $M_{0}, M_{1} > 0$,对 $X$ 上的简单函数 $f(x)$ 有
\begin{equation*}
\|Tf\|_{q_{j}} \leq M_{j}\|f\|_{p_{j}} \quad (j = 0, 1).
\end{equation*}
此时,对 $0 < \theta < 1$,$\|Tf\|_{q} \leq M_{0}^{1-\theta}M_{1}^{\theta}\|f\|_{p}$ 成立.其中 $p, q$ 由
\begin{equation*}
\frac{1}{p} = \frac{1 - \theta}{p_{0}} + \frac{\theta}{p_{1}}, \quad \frac{1}{q} = \frac{1 - \theta}{q_{0}} + \frac{\theta}{q_{1}}
\end{equation*}
给出.
\end{theorem}

下面的定理是 de Leeuw \cite{ch3:17}(cf. \cite{ch3:27})的结果.

\begin{theorem}\label{thm:3.5.17}
设 $1 < p < 2$.设 $m$ 是 $\mathbf{R}$ 上具有 $2\pi$ 周期的函数,$m \in L^{\infty}(\mathbf{R})$.此时,对任意 $\{a_{n}\} \in \ell^{p}$,令 $\hat{f}(x) = \sum a_{n}e^{inx}$,存在某个 $\{b_{n}\} \in \ell^{p}$ 使 $m(x)\hat{f}(x) = \sum_{n \in \mathbf{Z}}b_{n}e^{inx}$ 成立的充要条件是:对任意 $L^{p}(\mathbf{R})$ 的元 $g$ 的 Fourier 变换 $\hat{g}$,$m\hat{g}$ 是 $L^{p}(\mathbf{R})$ 的 Fourier 变换.其中,设 $\ell^{p} = \{a = \{a_{n}\}_{n=-\infty}^{\infty}: \|a\|_{\ell^{p}} = (\sum_{n \in \mathbf{Z}}|a_{n}|^{p})^{1/p} < \infty\}$,$\hat{a}(x) = \sum_{n}a_{n}e^{inx}$.
\end{theorem}

\begin{proof}
证明条件是必要的.对 $a = \{a_{n}\} \in \ell^{p}$,令 $\widehat{T_{m}}a(x) = m(x)\widehat{a}(x)$,由闭图像定理,$T_{m}$ 是从 $\ell^{p}$ 到 $\ell^{p}$ 的有界线性算子.把这样的 $T_{m}$ 的全体记为 $\mathfrak{M}_{p}(\mathbf{Z})$.又把 $T_{m}$ 的算子范数记为 $\|T_{m}\|_{\mathfrak{M}_{p}(\mathbf{Z})}$(或 $\|m\|_{M_{p}(\mathbf{T})}$).于是令 $m_{n} = \frac{1}{2\pi}\int_{-\pi}^{\pi}m(x)e^{-inx}\dd x$($n \in \mathbf{Z}$),进一步对 $f \in L^{p}(\mathbf{R})$ 令 $Kf(x) = \sum_{n}m_{n}f(x + n)$.这里,
\begin{equation*}
\infty > \int_{-\infty}^{\infty}|f(x)|^{p}\dd x = \int_{-1/2}^{1/2}\sum_{n \in \mathbf{Z}}|f(x + n)|^{p}\dd x
\end{equation*}
故 $\{f(x + n)\} \in \ell^{p}$(a.e. $x \in [-\frac{1}{2}, \frac{1}{2}]$).又由 $T_{m}$ 是 $\ell^{p}$ 上的有界算子以及 Banach 空间 $\ell^{p}$ 上的连续线性泛函的空间是 Banach 空间 $\ell^{p'}$($1/p + 1/p' = 1$),得 $\{m_{n}\} \in \ell^{p'}$.因此由 Hölder 不等式,
\begin{align*}
|Kf(x)| &\leq \sum_{n}|m_{n}||f(x + n)| \leq \left(\sum|m_{n}|^{p'}\right)^{1/p'}\left(\sum|f(x + n)|^{p}\right)^{1/p}\\
&< \infty \quad \left(\text{a.e. } x \in \left[-\frac{1}{2}, \frac{1}{2}\right]\right).
\end{align*}
又
\begin{equation*}
\begin{aligned}
\int_{-\infty}^{\infty}|Kf(x)|^{p}\dd x &= \int_{[-\frac{1}{2}, \frac{1}{2}]}\sum_{n \in \mathbf{Z}}|Kf(x + n)|^{p}\dd x\\
&= \int_{[-\frac{1}{2}, \frac{1}{2}]}\sum_{n \in \mathbf{Z}}\left|\sum_{k}m_{k}f(x + n + k)\right|^{p}\dd x.
\end{aligned}
\end{equation*}
这里利用 $T_{m}$ 是从 $\ell^{p}$ 到 $\ell^{p}$ 的有界线性算子,
\begin{align*}
&\int_{-\infty}^{\infty}|Kf(x)|^{p}\dd x \leq \|T_{m}\|_{\mathfrak{M}_{p}(\mathbf{Z})}^{p}\int_{[-\frac{1}{2}, \frac{1}{2}]}\left(\sum_{n}|f(x + n)|^{p}\right)^{p/p}\dd x\\
&= 2\pi\|T_{m}\|_{\mathfrak{M}_{p}(\mathbf{Z})}^{p}\|f\|_{L^{p}(\mathbf{R})}^{p}.
\end{align*}
因此 $K$ 是 $L^{p}(\mathbf{R})$ 上的有界线性算子.其次证明 $\widehat{Kf}(\xi) = m(\xi)\widehat{f}(\xi)$($f \in C_{c}^{\infty}(\mathbf{R})$).设 $j$ 为自然数,令 $s_{j}(x) = 2\pi j\Delta_{1}(jx)$($\Delta_{1}$ 是定理~\ref{thm:3.3.1} 的记号),$\sigma_{j}(x) = \sum_{n \in \mathbf{Z}}s_{j}(x + 2n\pi)$.此时,因 $m$ 是 $2\pi$ 周期函数,
\begin{align*}
\sigma_{j} * m(x) &= \frac{1}{2\pi}\int_{-\pi}^{\pi}\sigma_{j}(y)m(x - y)\dd y\\
&= \frac{1}{2\pi}\int_{-\pi}^{\pi}\sum_{n}s_{j}(y + 2n\pi)m(x - y)\dd y\\
&= \frac{1}{2\pi}\int_{-\infty}^{\infty}s_{j}(y)m(x - y)\dd y = s_{j} * m(x).
\end{align*}
此时,不妨设 $s_{j} * m(x) \to m(x)$ a.e..事实上,对 $0 < a < b < \infty$,因 $m \in L^{1}([a, b])$,
\begin{equation*}
\int_{-\infty}^{\infty}|m(x - y) - m(x)|\chi_{[a, b]}(x)\dd x \to 0 \quad (y \to 0).
\end{equation*}
因此
\begin{align*}
\int_{-\infty}^{\infty}\left|s_{j} * m(x)\chi_{[a, b]}(x) - m(x)\chi_{[a, b]}(x)\right|\dd x&\leq \int_{-\infty}^{\infty}\int_{-\infty}^{\infty}s_{j}(y)|m(x - y) - m(x)|\chi_{[a, b]}(x)\dd y\dd x\\
&= \int_{-\infty}^{\infty}s_{j}(y)\left(\int_{-\infty}^{\infty}|m(x - y) - m(x)|\chi_{[a, b]}(x)\dd x\right)\dd y \to 0 \quad (j \to \infty).
\end{align*}
从而,存在子列 $\{j_{k}\}$ 使得 $s_{j_{k}} * m(x) \to m(x)$ a.e..这里利用 $\|\sigma_{j}\|_{L^{1}(\mathbf{T})} = 1$ 计算,由
\begin{equation*}
\|T_{\sigma_{j} * m}\|_{\mathfrak{M}_{p}(\mathbf{Z})} = \|\sigma_{j} * m\|_{M_{p}(\mathbf{T})} \leq \|\sigma_{j}\|_{L^{1}(\mathbf{T})}\|m\|_{M_{p}(\mathbf{T})} = \|m\|_{M_{p}(\mathbf{T})},
\end{equation*}
注意 $T_{\sigma_{j} * m}$ 是从 $\ell^{p}$ 到 $\ell^{p}$ 的有界线性算子.由此,设 $s_{j, n} = \widehat{\sigma}_{j}(n)$,
\begin{equation*}
K_{j}f(x) = \sum_{n}s_{j, n}m_{n}f(x + n),
\end{equation*}
则 $K_{j}$ 与 $K$ 一样是从 $L^{p}(\mathbf{R})$ 到 $L^{p}(\mathbf{R})$ 的有界线性算子.此时
\begin{equation*}
|K_{j}f(x) - Kf(x)| \leq \sum_{n}|s_{j, n} - 1||m_{n}f(x + n)|.
\end{equation*}
这里,由
\begin{equation*}
s_{j, n} = \left(\sin\frac{n}{2j}\right)^{2}\left(\frac{n}{2j}\right)^{-2}, \quad \sum_{n}|m_{n}f(x + n)| < \infty \ \text{ a.e. } x,
\end{equation*}
$K_{j}f(x) \to Kf(x)$ a.e.($j \to \infty$).又
\begin{align*}
\widehat{K_{j}f}(x) &= \frac{1}{2\pi}\int_{-\infty}^{\infty}\sum_{n}s_{j, n}m_{n}f(y + n)e^{-ixy}\dd y\\
&= \sum_{n}s_{j, n}m_{n}e^{-inx}\widehat{f}(x),
\end{align*}
利用 $\widehat{s}_{j}(n) = s_{j, n}$,
\begin{equation*}
\widehat{K_{j}f}(x) = (s_{j} * m)(x)\widehat{f}(x).
\end{equation*}
因此,由 Parseval 等式,
\begin{align*}
&\|K_{j_{k}}f - T_{m}f\|_{L^{2}(\mathbf{R})} = \|\widehat{K_{j_{k}}f} - m\widehat{f}\|_{L^{2}(\mathbf{R})}\\
&\qquad = \|s_{j_{k}} * m(x)\widehat{f}(x) - m(x)\widehat{f}(x)\|_{L^{2}(\mathbf{R})}^{2}\\
&\qquad = \|(s_{j_{k}} * m - m)\widehat{f}\|_{L^{2}(\mathbf{R})}^{2} \to 0 \quad (k \to \infty),
\end{align*}
得 $\widehat{Kf} = m\widehat{f}$ a.e..

其次证明逆命题.设满足 $\widehat{Kf} = m\widehat{f}$ 的算子 $K$ 是从 $L^{p}(\mathbf{R})$ 到 $L^{p}(\mathbf{R})$ 的有界线性算子.此时,把算子范数记为 $\|K\|_{\mathfrak{M}_{p}(\mathbf{R})}$ 或 $\|m\|_{M_{p}(\mathbf{R})}$.又对 $c = \{c_{n}\}$,$d = \{d_{n}\} \in \ell^{p}$(于是 $c_{n}$,$d_{n}$ 除有限个外为 0),令 $f(x) = \sum_{m}c_{m}\chi_{(-\frac{1}{2}, \frac{1}{2})}(x + m)$,$g(x) = \sum_{n}d_{n}\chi_{(-\frac{1}{2}, \frac{1}{2})}(x + n)$,则 $\|f\|_{L^{p}(\mathbf{R})}^{p} = \frac{1}{2\pi}\|c\|_{\ell^{p}}^{p}$,$\|g\|_{L^{p'}(\mathbf{R})}^{p'} = \frac{1}{2\pi}\|d\|_{\ell^{p'}}^{p'}$.进一步,由 Parseval 等式,
\begin{align*}
\left|\int_{-\infty}^{\infty}\widehat{Kf}(x)\overline{\widehat{g}(x)}\dd x\right| &= \left|\frac{1}{2\pi}\int Kf(y)\overline{g(y)}\dd y\right|\\
&\leq \|m\|_{M_{p}(\mathbf{R})}\|f\|_{p}\|g\|_{p'}\\
&= \frac{1}{2\pi}\|m\|_{M_{p}(\mathbf{R})}\|c\|_{\ell^{p}}\|d\|_{\ell^{p'}}.
\end{align*}
另一方面,
\begin{align*}
\int_{-\infty}^{\infty}\widehat{Kf}(x)\overline{\widehat{g}(x)}\dd x&= \int_{-\infty}^{\infty}m(x)\widehat{f}(x)\overline{\widehat{g}(x)}\dd x\\
&= \int_{-\infty}^{\infty}m(x)\left(\sum_{k}c_{k}e^{-ikx}\widehat{\chi_{(-\frac{1}{2}, \frac{1}{2})}}(x)\right)\overline{\left(\sum_{n}d_{n}e^{-inx}\widehat{\chi_{(-\frac{1}{2}, \frac{1}{2})}}(x)\right)}\dd x\\
&= 2\pi\sum_{k, n}c_{k}\overline{d_{n}}\int_{0}^{2\pi}m(x)\sum_{j}\left|\widehat{\chi_{(-\frac{1}{2}, \frac{1}{2})}}(x + 2j\pi)\right|^{2}e^{-i(n-k)x}\dd x\\
&= 2\pi\sum_{k, n}c_{k}\overline{d_{n}}m_{n - k}
\end{align*}
成立(利用 Parseval 等式,$\sum_{j}|\widehat{\chi_{(-\frac{1}{2}, \frac{1}{2})}}(x + 2j\pi)|^{2} = \frac{1}{4\pi^{2}}$).从而
\begin{equation*}
\sup_{\|d\|_{\ell^{p'}} \leq 1}\left|\sum_{k, n}c_{k}\overline{d_{n}}m_{n - k}\right| \leq \|m\|_{M_{p}(\mathbf{R})}\|c\|_{\ell^{p}},
\end{equation*}
故满足 $\widehat{T_{m}f} = m\widehat{f}$ 的算子 $T_{m}$ 是 $\ell^{p}$ 上的有界线性算子,且 $\|T_{m}\|_{\mathfrak{M}_{p}(\mathbf{Z})} = \|m\|_{M_{p}(\mathbf{T})} \leq \|m\|_{M_{p}(\mathbf{R})}$.
\end{proof}

关于 $M_{p}(\mathbf{Z})$ 的元,下面的事实成立.该定理出自 K. de Leeuw \cite{ch3:17},Igari \cite{ch3:25},但这里按 Wozniakowski \cite{ch3:56},Kaneko and Sato \cite{ch3:28} 给出证明.

\begin{theorem}\label{thm:3.5.18}
设 $1 \leq p < \infty$.设 $\mathbf{R}$ 上的有界函数 $\lambda$ 在每个点连续.此时,令 $\widehat{Tf}(\xi) = \lambda(\xi)\widehat{f}(\xi)$,$T$ 是从 $L^{p}(\mathbf{R})$ 到 $L^{p}(\mathbf{R})$ 的有界线性算子的充要条件是:对任意 $\varepsilon > 0$,由 $\widehat{T_{\varepsilon}}f(n) = \lambda(\varepsilon n)\widehat{f}(n)$ 定义的 $\widetilde{T_{\varepsilon}}$ 是 $L^{p}(\mathbf{T})$ 上的有界线性算子,且令 $\lambda_{\varepsilon}(n) = \lambda(\varepsilon n)$,关于其算子范数 $\|T_{\varepsilon}\|_{\mathfrak{M}_{p}(\mathbf{T})} = \|\lambda_{\varepsilon}\|_{M_{p}(\mathbf{Z})}$,$\sup_{\varepsilon > 0}\|\lambda_{\varepsilon}\|_{M_{p}(\mathbf{Z})} < \infty$ 成立.此时,$\|T\|_{M_{p}(\mathbf{R})} = \sup_{\varepsilon > 0}\|\lambda_{\varepsilon}\|_{M_{p}(\mathbf{Z})}$ 成立.其中,$\|T\|_{\mathfrak{M}_{p}(\mathbf{R})}$ 表示 $T$ 在 $L^{p}(\mathbf{R})$ 上的算子范数.
\end{theorem}

\begin{note}\label{note:3.5.19}
把关于 $\lambda$ 的条件``在每个点连续''放宽为``在 $\mathbf{Z}$ 上的点连续'',通过审视证明也可证明 $\lambda|_{\mathbf{Z}} \in M_{p}(\mathbf{Z})$.
\end{note}

\begin{proof}
设 $\lambda$ 是 $\mathbf{R}$ 上的有界连续函数,对任意 $\varepsilon > 0$,令 $\widehat{T_{\varepsilon}f}(\xi) = \lambda(\varepsilon\xi)\widehat{f}(\xi)$($f \in C_c(\mathbf{R})$),设 $T = T_1$ 是 $L^p(\mathbf{R})$ 上的有界线性算子.又由 $\widehat{T_{\varepsilon}g}(n) = \lambda(\varepsilon n)\widehat{g}(n)$($g \in C(\mathbf{T})$)定义算子 $\widetilde{T_{\varepsilon}}$.设 $F$ 为 $\mathbf{T}$ 上的三角多项式.取 $\chi_k \in C^\infty(\mathbf{R})$,$\chi_k \geq 0$,满足 $\chi_k(\tau) = 1$($|\tau| \leq 2k\pi$),$\chi_k(\tau) = 0$($|\tau| > 2(k+1)\pi$),令 $\chi_k^\eta(x) = \chi_k(\eta x)$($\eta > 0$).进一步令 $\gamma_{\eta,k}(x) = \chi_k^\eta(x)(\widetilde{T_{\varepsilon}}F)(x) - T_{\varepsilon}(\chi_k^\eta F)(x)$,并令
\begin{equation*}
\widetilde{\gamma_{\eta,k}}(\xi) = \frac{1}{2\pi}\int_{-\infty}^{\infty}\chi_k^\eta(x)(\widetilde{T_{\varepsilon}}F)(x)e^{i\xi x}\dd x - \frac{1}{2\pi}\int_{-\infty}^{\infty}T_{\varepsilon}(\chi_k^\eta F)(x)e^{i\xi x}\dd x = I - II.
\end{equation*}

于是计算 I,II.
\begin{align*}
\mathrm{I} &= \frac{1}{2\pi}\int_{-\infty}^{\infty}\chi_k^\eta(x)\sum_j\lambda(j\varepsilon)\widehat{F}(j)e^{ijx}e^{i\xi x}\dd x\\
&= \frac{1}{2\pi}\sum_j\lambda(j\varepsilon)\widehat{F}(j)\int_{-\infty}^{\infty}\chi_k^\eta(x)e^{i(j+\xi)x}\dd x\\
&= \sum_j\lambda(j\varepsilon)\widehat{F}(j)\widehat{\chi_k^\eta}(-j-\xi),
\end{align*}
其中,
\begin{equation*}
\widehat{\chi_k^\eta}(-j-\xi) = \frac{1}{2\pi}\int_{-\infty}^{\infty}\chi_k(x)e^{i(j+\xi)x/\eta}\frac{1}{\eta}\dd x,
\end{equation*}
所以
\begin{equation*}
\mathrm{I} = \sum_j\lambda(j\varepsilon)\widehat{F}(j)\widehat{\chi_k}\left(\frac{-j-\xi}{\eta}\right)\frac{1}{\eta}.
\end{equation*}
又,
\begin{align*}
\mathrm{II} &= \frac{1}{2\pi}\int_{-\infty}^{\infty}T_{\varepsilon}(\chi_k^\eta F)(x)e^{i\xi x}\dd x\\
&= \lambda(-\varepsilon\xi)\frac{1}{2\pi}\int_{-\infty}^{\infty}\chi_k^\eta(x)F(x)e^{i\xi x}\dd x\\
&= \lambda(-\varepsilon\xi)\sum_j\widehat{F}(j)\frac{1}{2\pi}\int_{-\infty}^{\infty}\chi_k^\eta(x)e^{i(j+\xi)x}\dd x
\end{align*}
\begin{equation*}
= \lambda(-\varepsilon\xi)\sum_j\widehat{F}(j)\eta^{-1}\widehat{\chi_k}\left(\frac{-j-\xi}{\eta}\right),
\end{equation*}
所以,最终
\begin{equation*}
\mathrm{I} - \mathrm{II} = \sum_j\widehat{F}(j)\eta^{-1}\widehat{\chi_k}\left(\frac{-j-\xi}{\eta}\right)(\lambda(j\varepsilon) - \lambda(-\xi\varepsilon)).
\end{equation*}
因此,
\begin{equation*}
\widetilde{\gamma_{\eta,k}}(\xi) = \sum_j\widehat{F}(j)\eta^{-1}\widehat{\chi_k}\left(\frac{\xi - j}{\eta}\right)(\lambda(j\varepsilon) - \lambda(\xi\varepsilon)),
\end{equation*}
所以
\begin{align*}
\left\|\widetilde{\gamma_{\eta,k}}\right\|_{1} &\leq \sum\left|\widehat{F}(j)\right|\int_{-\infty}^{\infty}\eta^{-1}\left|\widehat{\chi_k}\left(\frac{\xi - j}{\eta}\right)\right||\lambda(j\varepsilon) - \lambda(\xi\varepsilon)|\dd\xi\\
&\leq \sum\left|\widehat{F}(j)\right|\int_{-\infty}^{\infty}|\widehat{\chi_k}(\xi)||\lambda(j\varepsilon) - \lambda(j\varepsilon + \eta\xi\varepsilon)|\dd\xi.
\end{align*}
因此
\begin{equation*}
|\gamma_{\eta,k}(x)| \leq \|\widetilde{\gamma_{\eta,k}}\|_{1} \leq \sum_j\left|\widehat{F}(j)\right|\int_{-\infty}^{\infty}|\widehat{\chi_k}(\xi)||\lambda(j\varepsilon) - \lambda(j\varepsilon + \eta\xi\varepsilon)|\dd\xi.
\end{equation*}
因 $F$ 是三角多项式,固定 $k$ 时,若 $\eta \to 0$,则 $\|\widetilde{\gamma_{\eta,k}}\|_{1} \to 0$.于是,对 $a > 0$,取 $\eta > 0$ 使 $\|\widetilde{\gamma_{\eta,k}}\|_{1} < a$,设 $t > a$,则
\begin{equation*}
\left|\{x: |T_{\varepsilon}(\chi_k^\eta F)(x)| > t - \|\widetilde{\gamma_{\eta,k}}\|_{1}\}\right| \geq \left|\{x: |\chi_k^\eta(\widetilde{T_{\varepsilon}}F)(x)| > t\}\right|
\end{equation*}
(对集合 $E$,$|E|$ 表示 $E$ 的 Lebesgue 测度).从而
\begin{align*}
\mathrm{III} &= \int_{a}^{\infty}p\left|\left\{\left|\chi_k^\eta(\widetilde{T_{\varepsilon}}F)(x)\right| > t\right\}\right|t^{p-1}\dd t\\
&\leq \int_{a}^{\infty}p\left|\{|T_{\varepsilon}(\chi_k^\eta F)(x)| > t - \|\widetilde{\gamma_{\eta,k}}\|_{1}\}\right|t^{p-1}\dd t\\
&\leq \int_{a - \|\widetilde{\gamma_{\eta,k}}\|_{1}}^{\infty}p\left|\{|T_{\varepsilon}(\chi_k^\eta F)(x)| > t\}\right|(t + \|\widetilde{\gamma_{\eta,k}}\|_{1})^{p-1}\dd t
\end{align*}
\begin{align*}
&= \int_{a - \|\widetilde{\gamma_{\eta,k}}\|_{1}}^{\infty}p\left(t|\{|T_{\varepsilon}(\chi_k^\eta F)(x)| > t\}|^{1/p}\right)^{p-1}\left(1 + \frac{\|\widetilde{\gamma_{\eta,k}}\|_{1}}{t}\right)^{p-1}\\
&\leq \left(1 + \frac{\|\widetilde{\gamma_{\eta,k}}\|_{1}}{a - \|\widetilde{\gamma_{\eta,k}}\|_{1}}\right)^{p-1}\|T_{\varepsilon}(\chi_k^\eta F)\|_{p}^{p}.
\end{align*}
这里,因 $T_{\varepsilon}(\chi_{k}^{\eta}F)(x) = (T(\chi_{k}^{\eta}F)^{\varepsilon})(x/\varepsilon)$(其中 $(\chi_{k}^{\eta}F)^{\varepsilon}(x) = (\chi_{k}^{\eta}F)(\varepsilon x)$),
\begin{align*}
&\mathrm{III} \leq \left(1 + \frac{\|\widetilde{\gamma_{\eta,k}}\|_{1}}{a - \|\widetilde{\gamma_{\eta,k}}\|_{1}}\right)^{p}\left\|(T(\chi_{k}^{\eta}F)^{\varepsilon})\left(\frac{\cdot}{\varepsilon}\right)\right\|_{p}^{p}\\
&\quad \leq \left(1 + \frac{\|\widetilde{\gamma_{\eta,k}}\|_{1}}{a - \|\widetilde{\gamma_{\eta,k}}\|_{1}}\right)^{p}\varepsilon\|T\|_{\mathfrak{M}_{p}(\mathbf{R})}^{p}\|(\chi_{k}^{\eta}F)^{\varepsilon}\|_{p}^{p}\\
&\quad = \left(1 + \frac{\|\widetilde{\gamma_{\eta,k}}\|_{1}}{a - \|\widetilde{\gamma_{\eta,k}}\|_{1}}\right)^{p}\|T\|_{\mathfrak{M}_{p}(\mathbf{R})}^{p}\|\chi_{k}^{\eta}F\|_{p}^{p}.
\end{align*}
这里,由 $0 \leq \chi_{k}^{\eta} \leq 1$,$\supp\chi_{k}^{\eta} \subset \left[\frac{-2(k+1)\pi}{\eta}, \frac{2(k+1)\pi}{\eta}\right]$,对 $Q_{l} = [-\pi, \pi] + 2l\pi$($l \in Z$),满足 $Q_{l} \cap \left[\frac{-2(k+1)\pi}{\eta}, \frac{2(k+1)\pi}{\eta}\right] \neq \emptyset$ 的 $l$ 的个数至多为 $\left[2\left(\frac{k+1}{\eta}\right)\right]$.因此
\begin{equation*}
\|\chi_{k}^{\eta}F\|_{L^{p}(\mathbf{R})} \leq \left\{2\left(\frac{k+1}{\eta} + 1\right)\right\}^{1/p}\|F\|_{L^{p}(\mathbf{T})}.
\end{equation*}
另一方面,由 $\chi_{k}^{\eta}(x) = 1$($x \in \left[\frac{-2k\pi}{\eta}, \frac{2k\pi}{\eta}\right]$),满足 $Q_{j} \subset \left[\frac{-2k\pi}{\eta}, \frac{2k\pi}{\eta}\right]$ 的 $Q_{j}$ 的个数至少为 $\left[2\left(\frac{k}{\eta} - 1\right)\right]$,所以
\begin{align*}
\left|\left\{\left|\chi_{k}^{\eta}(\widetilde{T_{\varepsilon}}F)\right| > t\right\}\right| &\geq \left|\bigcup_{j}\left\{Q_{j} \cap \left\{\left|\widetilde{T_{\varepsilon}}F\right| > t\right\}\right\}: Q_{j} \subset \left[\frac{-2k\pi}{\eta}, \frac{2k\pi}{\eta}\right]\right|\\
&\geq \left\{2\left(\frac{k}{\eta} - 1\right)\right\}\left|\left\{\left|\widetilde{T_{\varepsilon}}F\right| > t\right\}\right|.
\end{align*}
综合以上,
\begin{align*}
&\left\{2\left(\frac{k}{\eta} - 1\right)\right\}^{1/p}\left\{\int_{a}^{\infty}p\left|\left\{\left|\widetilde{T_{\varepsilon}}F\right| > t\right\}\right|^{1/p}t^{p-1}\dd t\right\}^{1/p}\\
&\qquad \leq \left\{2\left(\frac{k+1}{\eta} + 1\right)\right\}^{1/p}\left\{1 + \frac{\|\widetilde{\gamma_{\eta,k}}\|_{1}}{a - \|\widetilde{\gamma_{\eta,k}}\|_{1}}\right\}\|T\|_{\mathfrak{M}_{p}(\mathbf{R})}\|F\|_{p}.
\end{align*}
因此
\begin{align*}
\{2(k - \eta)\}^{1/p}\left\{\int_{a}^{\infty}p\left|\left\{\left|\widetilde{T_{\varepsilon}}F\right| > t\right\}\right|^{1/p}t^{p-1}\dd t\right\}^{1/p}&\leq \{2(k + 1 + \eta)\}^{1/p}\left\{1 + \frac{\|\widetilde{\gamma_{\eta,k}}\|_{1}}{a - \|\widetilde{\gamma_{\eta,k}}\|_{1}}\right\}\|T\|_{\mathfrak{M}_{p}(\mathbf{R})}\|F\|_{p}.
\end{align*}
这里令 $\eta \to 0$,得
\begin{align*}
(2k)^{1/p}\left\{\int_{a}^{\infty}p\left|\left\{\left|\widetilde{T_{\varepsilon}}F\right| > t\right\}\right|^{1/p}t^{p-1}\dd t\right\}^{1/p}&\leq \{2(k + 1)\}^{1/p}\|T\|_{\mathfrak{M}_{p}(\mathbf{R})}\|F\|_{p}.
\end{align*}
除以 $k$ 后,令 $k \to \infty$,再在 $a \to 0$ 时利用 Fatou 引理,得
\begin{equation*}
\|\widetilde{T_{\varepsilon}}F\|_{p} \leq \|T\|_{\mathfrak{M}_{p}(\mathbf{R})}\|F\|_{p}.
\end{equation*}

反之,设 $\sup_{\varepsilon>0}\|\widetilde{T_{\varepsilon}}\|_{M_{p}(\mathbf{T})}<\infty$.取 $f\in C_{c}^{\infty}(\mathbf{R})$,令 $F_{\varepsilon}(x)=\varepsilon^{-1}\sum_{k\in\mathbf{Z}}f(\varepsilon^{-1}(x+2k\pi))$,则 $F_{\varepsilon}\in L^{1}(\mathbf{T})$.于是,
\begin{align*}
\widehat{F_{\varepsilon}}(j) &= \frac{1}{2\pi}\int_{-\pi}^{\pi}F_{\varepsilon}(x)e^{-ijx}\dd x\\
&= \frac{1}{2\pi}\sum_{k}\varepsilon^{-1}\int_{-\pi}^{\pi}f(\varepsilon^{-1}(x+2k\pi))e^{-ijx}\dd x\\
&= \frac{1}{2\pi}\int_{-\infty}^{\infty}\varepsilon^{-1}f(\varepsilon^{-1}x)e^{-ijx}\dd x = \widehat{f}(\varepsilon j),
\end{align*}
得 $\widehat{F_{\varepsilon}}(j)=\widehat{f}(\varepsilon j)$.因此 $F_{\varepsilon}(x)=\sum_{j}\widehat{f}(\varepsilon j)e^{ijx}$.这里,考虑 Riemann 积分的定义,
\begin{align*}
|Tf(x)| &= \left|\int_{-\infty}^{\infty}\lambda(\xi)\widehat{f}(\xi)e^{ix\xi}\dd\xi\right|\\
&= \lim_{\varepsilon \to 0}\left|\sum_{j}\lambda(\varepsilon j)\widehat{f}(\varepsilon j)e^{ix\varepsilon j}\varepsilon\right|
\end{align*}
\begin{align*}
&= \lim_{\varepsilon \to 0}\left|\sum_{j}\lambda(\varepsilon j)\widehat{F_{\varepsilon}}(j)e^{i\varepsilon xj}\varepsilon\right|\\
&= \lim_{\varepsilon \to 0}\left|\widetilde{T_{\varepsilon}}(\varepsilon F_{\varepsilon})(\varepsilon x)\right|.
\end{align*}
进一步,由 $\lim_{\varepsilon\to0}\chi_{[-\pi,\pi)}(\varepsilon x)=1$($x\in\mathbf{R}$),
\begin{equation*}
|Tf(x)| = \lim_{\varepsilon \to 0}\left|\widetilde{T_{\varepsilon}}(\varepsilon F_{\varepsilon})(\varepsilon x)\right|\chi_{[-\pi,\pi)}(\varepsilon x),
\end{equation*}
所以
\begin{align*}
\|Tf\|_{L^{p}(\mathbf{R})} &= \left(\int_{-\infty}^{\infty}|Tf|^{p}\frac{\dd x}{2\pi}\right)^{1/p}\\
&\leq \liminf_{\varepsilon \to 0}\left(\frac{1}{2\pi}\int_{-\pi}^{\pi}\left|\widetilde{T_{\varepsilon}}(\varepsilon F_{\varepsilon})(x)\right|^{p}\varepsilon^{-1}\dd x\right)^{1/p}\\
&\leq \liminf_{\varepsilon \to 0}\varepsilon^{-1/p}\|\widetilde{T_{\varepsilon}}(\varepsilon F_{\varepsilon})\|_{L^{p}(\mathbf{T})}\\
&\leq \left(\sup_{\varepsilon}\|\widetilde{T_{\varepsilon}}\|_{\mathfrak{M}_{p}(\mathbf{T})}\right)\liminf_{\varepsilon \to 0}\varepsilon^{-1/p}\|\varepsilon F_{\varepsilon}\|_{L^{p}(\mathbf{T})}\\
&\leq \left(\sup_{\varepsilon}\|\widetilde{T_{\varepsilon}}\|_{\mathfrak{M}_{p}(\mathbf{T})}\right)\|f\|_{L^{p}(\mathbf{R})}.
\end{align*}
因此
\begin{equation*}
\|Tf\|_{L^{p}(\mathbf{R})} \leq \left(\sup_{\varepsilon > 0}\|\widetilde{T_{\varepsilon}}\|_{\mathfrak{M}_{p}(\mathbf{T})}\right)\|f\|_{L^{p}(\mathbf{R})}.
\end{equation*}
\end{proof}

下面研究乘子空间 $M_{p}(\mathbf{Z})$($1 < p < 2$)的作用函数.为此,先叙述若干引理(\cite{ch3:24},cf. \cite{ch3:43}, \cite{ch3:48}).

\begin{lemma}\label{lem:3.5.20}
设 $1 < p < 2$.存在 $K_{p}(>1)$,使得对每个自然数 $j$,存在满足下列条件的 $\mathbf{T}$ 上的三角多项式 $\phi = \phi_{p,j}$:

(i) $\widehat{\phi}(n) \geq 0$($n \in \mathbf{Z}$),$\|\phi\|_{L^1(\mathbf{T})} = 1$.

(ii) $\|\exp(ij\phi)\|_{M_p(\mathbf{T})} > K_p^j$.

这里,$\exp(ij\phi) \in M(\mathbf{T})$ 由 $\exp(ij\phi)^{\wedge}(n) = \exp(ij\widehat{\phi}(n))$($n \in \mathbf{Z}$)定义.
\end{lemma}

\begin{proof}
分几个阶段叙述.

第 1 阶段.令 $g(x)=e^{i\cos x}$,因 $g\in A(\mathbf{T})$,
\begin{equation*}
K_{p}^{\prime} = \left(\sum|\widehat{g}(n)|^{p}\right)^{1/p} < \sum|\widehat{g}(n)| < \infty.
\end{equation*}
又由 $1 < p < 2$,$K_{p}^{\prime} > \left(\sum|g(n)|^{2}\right)^{1/2} = \|g\|_{2} = 1$,故 $K_{p}^{\prime} > 1$.现在取 $0 < \mu < 1$ 使 $K_{p} = \mu K_{p}^{\prime} > 1$.

第 2 阶段.令 $Q(x) = \sum_{j=1}^{r}\cos n_{j}x$,适当选取 $n_{1}, \ldots, n_{r}$,令 $g = e^{iQ}$,则 $A_{p}(g) = \left(\sum|\widehat{g}(n)|^{p}\right)^{1/p} > K_{p}^{r}$($r = 1, 2, \ldots$).事实上,令 $n_{1} = 1$,由第 1 阶段即成立.设 $n = k < r$ 时成立.此时,对 $f = e^{i\sum_{j=1}^{k}\cos n_{j}x}$,$g = e^{i\cos n_{k+1}x}$,由 $A_{p}(fg) \leq \|f\|_{A(\mathbf{T})}A_{p}(g)$,$f$ 不妨取为三角多项式.此时,把 $n_{k+1}$ 取得足够大,$A_{p}(fg) = A_{p}(f)A_{p}(g)$ 成立,故可做到 $A_{p}(fg) > K_{p}^{k}\mu A_{p}(g) > K_{p}^{k+1}$.从而 $A_{p}\left(e^{i\sum_{j=1}^{k+1}\cos n_{j}x}\right) > K_{p}^{k+1}$.

第 3 阶段.对上述 $K_{p} > 1$,存在 $\mathbf{T}$ 上的概率测度 $\mu$ 满足 $\widehat{\mu}(n) \in \mathbf{R}$($n \in \mathbf{Z}$)且 $\|\exp(ij\mu/2)\|_{\mathfrak{M}_{p}(\mathbf{T})} > K_{p}^{j}$.事实上,由第 2 阶段,存在 $\mathbf{T}$ 上的三角多项式 $Q$ 满足 $A_{p}(e^{iQ}) > K_{p}^{j}$($\|Q\|_{A(\mathbf{T})} \leq j$).把 $T_{\varphi}$ 看作由 $\widehat{T_{\varphi}}f(x) = \varphi(x)\widehat{f}(x)$($\widehat{f}(x) = \sum f_n e^{inx}$)给出的 $\ell^p$ 上的算子.其算子范数为 $\|T_{\varphi}\|_{\mathfrak{M}_{p}(\mathbf{Z})} = \|\varphi\|_{M_p(\mathbf{T})} \geq A_p(\varphi)$,且 $\|Q\|_{A(\mathbf{T})} = \|Q\|_{M_1(\mathbf{T})}$(cf. \cite{ch3:9},引理~\ref{lem:3.2.11}),故令 $\varphi = \frac{1}{j}Q$,得 $\|e^{ij\varphi}\|_{M_p(\mathbf{T})} > K_p^j$($j = 1, 2, \ldots$).此时,对该 $\varphi$ 应用定理~\ref{thm:3.5.17},再对其应用定理~\ref{thm:3.5.18},适当取 $\varepsilon > 0$,则存在概率测度 $\mu$ 满足 $\widehat{\mu}(n) = \varphi(\varepsilon n)$,$\widehat{\mu}(n) \in \mathbf{R}$($n \in \mathbf{Z}$)且 $\|e^{ij\mu}\|_{\mathfrak{M}_{p}(\mathbf{T})} > K_p^j$($j = 1, 2, \ldots$).进一步,
\begin{equation*}
\|e^{ij\mu/2}\|_{\mathfrak{M}_{p}(\mathbf{T})}^{2} \geq \|e^{ij\mu}\|_{\mathfrak{M}_{p}(\mathbf{T})} > K_{p}^{j},
\end{equation*}
所以 $\|e^{ij\mu/2}\|_{\mathfrak{M}_{p}(\mathbf{T})} > \left(\sqrt{K_{p}}\right)^{j}$.这里令 $\nu = (\delta_{0} + \mu)/2$,则 $\widehat{\nu}(n) \geq 0$($n \in \mathbf{Z}$),$\nu$ 是概率测度,且 $\|\exp(ij\nu)\|_{\mathfrak{M}_{p}(\mathbf{T})} > \sqrt{K_{p}}^{j}$($j = 1, 2, \ldots$).从而,取次数足够高的 Fejér 核 $\sigma_{n}$,令 $\phi = \nu * \sigma_{n}$,把 $\sqrt{K_{p}}$ 重新记为 $K_{p}$,即得结果.
\end{proof}

\begin{theorem}\label{thm:3.5.21}
设 $1 < p < 2$,则存在 $\mathbf{T}$ 上的概率测度 $\mu$($\hat{\mu} \geq 0$)和常数 $K_{p} > 1$,使得令
\begin{equation*}
Tf = e^{i\mu} * f \quad (f \in L^{p}(\mathbf{T})),
\end{equation*}
算子 $T$ 在 $\mathfrak{M}_{p}(\mathbf{T})$ 中的谱半径 $\rho(T)$ 满足 $\rho(T) \geq K_{p}$.
\end{theorem}

\begin{proof}
设 $Q$ 为有理数的集合.设 $\{p_{n}\}$ 为使任意 $p \in \mathbf{Q} \cap (1, 2)$ 出现无穷多次的数列.由引理~\ref{lem:3.5.20},对自然数 $j$ 存在 $\mathbf{T}$ 上的三角多项式 $\phi_{j}$ 使得 $\widehat{\phi_{j}} \geq 0$,$\|\phi_{j}\|_{L^{1}(\mathbf{T})} = 1$,$\|\exp(ij\phi_{j})\|_{M_{p}(\mathbf{T})} > K_{p}^{j}$($p = p_{j}$).于是,对 $p = p_{j}$,取三角多项式 $f_{j}$ 使得
\begin{equation*}
\|f_{j}\|_{p} = 1, \quad \left\|\sum_{k}\exp(ij\widehat{\phi_{j}}(k))\widehat{f_{j}}(k)\exp(ikt)\right\|_{p} > K_{p}^{j}.
\end{equation*}
此时,存在自然数 $m_{j}$ 使得
\begin{equation*}
\left(\supp\widehat{\phi_{j}}\right) \cup \left(\supp\widehat{f_{j}}\right) \subset \{-m_{j}, -m_{j} + 1, \ldots, m_{j} - 1, m_{j}\}.
\end{equation*}
现在,设 $\{r_{j}\} \subset \mathbb{N}$ 为急剧增大的序列,则满足
\begin{equation*}
d\mu_{n}(t) = \phi_{1}(r_{1}t)\phi_{2}(r_{2}t)\cdots\phi_{n}(r_{n}t)\frac{\dd t}{2\pi}
\end{equation*}
的 $\mu_{n}\in M(\mathbf{T})$ 关于 $*$-弱拓扑收敛于概率测度 $\mu\in M(\mathbf{T})$,且满足
\begin{equation*}
\widehat{\mu}(k_{1}r_{1} + \cdots + k_{n}r_{n}) = \prod_{j=1}^{n}\widehat{\phi_{j}}(k_{j}) \quad (|k_{j}| \leq m_{j},\ j = 1, 2, \ldots, n).
\end{equation*}
进一步,令 $g_{j}(t)=f_{j}(r_{j}t)$,
\begin{align*}
\left\|\sum_{k}\exp(ij\widehat{\mu}(k))\widehat{g_{j}}(k)\exp(ikt)\right\|_{p} &= \left\|\sum_{k}\exp(ij\widehat{\phi_{j}}(k))\widehat{f_{j}}(k)\exp(ikt)\right\|_{p}\\
&> K_{p}^{j}.
\end{align*}
因此 $\|\exp(ij\mu)\|_{M_{p}(\mathbf{T})}>K_{p}^{j}$($p=p_{j}$),由 Riesz--Thorin 插值定理,对所有 $p\in(1,2)$ 有 $\|\exp(ij\mu)\|_{M_{p}(\mathbf{T})}>K_{p}^{j}$.从而由谱半径的定义即得结果.
\end{proof}

\begin{note}\label{note:3.5.22}
对上述 $\mu$,由 Riesz--Thorin 插值定理,下式成立:
\begin{equation*}
\sup\left\{\|e^{ij\mu}\|_{\mathfrak{M}_{p}(\mathbf{T})}: \mu \in M(\mathbf{T}),\ \hat{\mu}(\mathbf{Z}) \subset \mathbf{R},\ \|\mu\|_{M(\mathbf{T})} \leq 1\right\} \leq e^{|j|(2/p - 1)} \quad (j \in \mathbf{Z}).
\end{equation*}
\end{note}

\begin{lemma}\label{lem:3.5.23}
在 $[-1, 1]$ 上定义的函数 $\Phi$ 若作用于 $M_{p}(\mathbf{Z})$,则连续.
\end{lemma}

\begin{proof}
设 $\Phi$ 在 $a \in [-1, 1]$ 处不连续.此时,存在 $\{a_j\} \subset [-1, 1]$ 和 $B \in \mathbf{R}$ 使得
\begin{equation*}
\Phi(a) \neq B, \quad a_i \neq a_j \ (i \neq j), \quad \sum_{j=-\infty}^{\infty}|a_j - a| < \infty, \quad \sum_{j=-\infty}^{\infty}|\Phi(a_j) - B| < \infty.
\end{equation*}
此时不妨设 $\Phi(a)=0$.又,适当选取 $\{\varepsilon_{j}\}_{j=-\infty}^{\infty}$($\varepsilon_{j}=\pm1$),取 $f \in L^{p}(\mathbf{T})$ 使得
\begin{equation*}
\sum_{j=-\infty}^{\infty}\varepsilon_{j}\widehat{f}(j)e^{ijx} \notin L^{p}(\mathbf{T})
\end{equation*}
($f$ 的存在性参见 \cite{ch3:24}).现在,令 $\eta_{j} = \Phi(a_{j})$($\varepsilon_{j} = 1$),$\eta_{j} = 0$($\varepsilon_{j} = -1$),则 $\sum_{-\infty}^{\infty}\widehat{f}(j)\eta_{j}e^{ijx} \notin L^{p}(\mathbf{T})$.事实上,由
\begin{align*}
\sum_{j=-\infty}^{\infty}\widehat{f}(j)\eta_{j}e^{ijx}&= \frac{B}{2}\sum_{j=-\infty}^{\infty}\widehat{f}(j)e^{ijx} + \frac{B}{2}\sum_{j=-\infty}^{\infty}\widehat{f}(j)\varepsilon_{j}e^{ijx} + \sum_{\varepsilon_{j}=1}\widehat{f}(j)(\Phi(a_{j}) - B)e^{ijx},
\end{align*}
右边第 1 项,第 3 项是 $L^{p}(\mathbf{T})$ 的元,而第 2 项不是.于是,令 $\varphi(j)=a_{j}$($\varepsilon_{j}=1$),$\varphi(j)=a$($\varepsilon_{j}\neq1$),则对 $g\in L^{1}(\mathbf{T})$,
\begin{equation*}
\sum_{j}\widehat{g}(j)\varphi(j)e^{ijx} = \sum_{j}\widehat{g}(j)(\varphi(j) - a)e^{ijx} + a\sum_{j}\widehat{g}(j)e^{ijx}.
\end{equation*}
因此右边是 $L^{1}(\mathbf{T})$ 的元,故 $\varphi \in M_{1}(\mathbf{Z})$.从而由 $\Phi$ 是 $M_{p}(\mathbf{Z})$ 的作用函数,$\Phi(\varphi) \in M_{p}(\mathbf{Z})$.可是,$\Phi(\varphi(j)) = \Phi(a_j) = \eta_j$,$\{\eta_j\} \notin M_p(\mathbf{Z})$,故 $\Phi$ 不作用于 $M_p(\mathbf{Z})$,矛盾.
\end{proof}

$M_{p}(\mathbf{Z})$ 的作用函数如下确定.

\begin{theorem}\label{thm:3.5.24}
$[-1, 1]$ 上的函数 $\Phi$ 作用于 $M_{p}(\mathbf{Z})$ 的充要条件是:与某个整函数一致.
\end{theorem}

\begin{proof}
设 $\Phi$ 是作用于 $[-1, 1]$ 上的 $M_p(\mathbf{Z})$ 的函数.考虑函数 $\Phi(\varepsilon\sin t)$($0 < \varepsilon < 1$),它是 $\mathbf{R}$ 上的 $2\pi$ 周期函数且作用于 $M_p(\mathbf{Z})$.因此,不妨从一开始假设 $\Phi$ 是 $\mathbf{R}$ 上的 $2\pi$ 周期函数,只需证明它与某个整函数一致(此后的讨论与 $M(\mathbf{T})$ 的作用函数的确定情形相同).此时,由引理~\ref{lem:3.5.23} 不妨设 $\Phi$ 连续.又,像引理~\ref{lem:3.5.13} 的证明那样考虑,对任意 $\lambda \in \mathbf{R}$ 存在 $C_\lambda > 0$ 使得 $\|\Phi(\lambda\hat{\mu} + a)\|_{M_p(\mathbf{Z})} \leq C_\lambda$($-\pi \leq a \leq \pi$).于是,把 $\Phi$ 的 Fourier 级数记为 $\Phi \sim \sum a_n e^{int}$.设 $j$ 为任意自然数,$n \in \mathbf{Z}$,则
\begin{align*}
a_{n}\exp(inj\widehat{\mu}) &= \frac{1}{2\pi}\int_{-\pi}^{\pi}\Phi(x)e^{-in(x-j\widehat{\mu})}\dd x\\
&= \frac{1}{2\pi}\int_{-\pi}^{\pi}\Phi(x + j\widehat{\mu})e^{-ijnx}\dd x.
\end{align*}
由前半部分,
\begin{equation*}
\|a_{n}\exp(inj\widehat{\mu})\|_{M_{p}(\mathbf{Z})} \leq \frac{1}{2\pi}\int\|\Phi(x + j\widehat{\mu})\|_{M_{p}(\mathbf{Z})}\dd x \leq C_{j}.
\end{equation*}
因此
\begin{equation*}
|a_{n}| \leq C_{j}\|\exp(inj\widehat{\mu})\|_{M_{p}(\mathbf{Z})}^{-1}.
\end{equation*}
这里,由引理~\ref{lem:3.5.20} 的证明,
\begin{equation*}
\|\exp(inj\widehat{\mu})\|_{M_{p}(\mathbf{Z})}^{-1} \leq C\exp(-|n|j\log K_{p}),
\end{equation*}
所以
\begin{equation*}
|a_{n}| \leq C_{j}'\exp(-|n|j\log K_{p}) \quad (n \in \mathbf{Z})
\end{equation*}
被证明.从而 $\Phi$ 在 $[-1, 1]$ 上与整函数一致.
\end{proof}

\subsection{\texorpdfstring{$M(p,q)$($1<p<q<\infty$)的作用函数}{M(p,q)(1<p<q<∞)的作用函数}}\label{sec:3.5.5}

设 $1 \leq p < q < \infty$.把从 $L^{p}(\mathbf{T})$ 到 $L^{q}(\mathbf{T})$ 的平移不变有界线性算子的全体记为 $\mathfrak{M}_{pq}$($= \{T \in \mathfrak{M}_{p}(\mathbf{T}): L^{p}(\mathbf{T})$ 到 $L^{q}(\mathbf{T})$ 的有界线性算子$\}$),它关于算子范数成为交换 Banach 代数,且 $\mathfrak{M}_{pq} \subset \mathfrak{M}_{p}(\mathbf{T})$.又定义 $M(p,q) = \{\hat{T} \in \ell^{\infty}: T \in \mathfrak{M}_{pq}(\mathbf{T})\}$($\subset M_{p}(\mathbf{Z})$).当 $p = 1$ 时,$\mathfrak{M}_{pq} = L^{q}(\mathbf{T})$(cf. \cite{ch3:9}).Rudin \cite{ch3:41},Rider \cite{ch3:40} 研究了 $M(1, q)$ 的作用函数.这里叙述推广了他们结果的 Igari--Sato \cite{ch3:26} 的结果.另外,$M(p,q)$ 的范数由 $\mathfrak{M}_{pq}$ 引入.对函数 $\phi$ 与算子 $T$,算子 $\phi \circ T$ 由 $\widehat{\phi \circ T}(n) = \phi(\hat{T}(n))$($n \in \mathbf{Z}$)定义,下文使用.

关于 $M(p,q)$ 的作用函数,下面的事实成立.

\begin{theorem}\label{thm:3.5.25}
(1) 设 $1 < p < q \leq 2$.设 $\beta_0 = (1/q - 1/2)/(1/p - 1/q)$,$n_0$ 为满足 $n_0 \geq \beta_0$ 的最小整数.此时,设 $\alpha_1, \ldots, \alpha_{n_0}$ 为任意复数,取 $[-1, 1]$ 上的有界复值函数 $\Phi_0(t)$,则
\begin{equation*}
\Phi(t) = \alpha_{1}t + \cdots + \alpha_{n_0}t^{n_0} + |t|^{\beta_0 + 1}\Phi_0(t)
\end{equation*}
作用于 $M(p,q)$.

(2) 设 $1 < p < 2 < q < \infty$.$[-1, 1]$ 上的复值函数 $\Phi$ 对复数 $\alpha$ 与
\begin{equation*}
\beta_{1} = \min\left\{\frac{\frac{1}{2} - \frac{1}{q}}{\frac{1}{p} - \frac{1}{2}}, \frac{\frac{1}{p} - \frac{1}{2}}{\frac{1}{2} - \frac{1}{q}}\right\}
\end{equation*}
以及有界函数 $\Phi_{0}$,若写成 $\Phi(t)=\alpha t+|t|^{\beta_{1}+1}\Phi_{0}(t)$,则 $\Phi$ 作用于 $M(p,q)$.
\end{theorem}

由此,关于 $M(p,q)$ 的元的谱,经若干计算,下面的事实成立(cf. \cite{ch3:26}).

\begin{corollary}\label{cor:3.5.26}
\begin{sloppypar}
设 $1 < p < q < \infty$.$M(p,q)$ 是不含单位元的交换 Banach 代数,对 $M(p,q) \ni f$,$\mathrm{Sp}(f, M(p,q)) = \overline{f(\mathbf{Z})}$.
\end{sloppypar}
\end{corollary}

\begin{proof}
由 Riesz--Thorin 插值定理和 Riesz 表示定理,不妨设 $q < 2$.设 $0 \neq \lambda \notin \widehat{T}(\mathbf{Z})$,$\lambda \in \mathrm{Sp}(T, \mathfrak{M}_{pq})$.此时,取足够大的自然数 $N$,由 Riesz--Thorin 插值定理,$T^{N} \in \mathfrak{M}_{p2}(\mathbf{T})$(cf. 引理~\ref{lem:3.5.29}).于是,令
\begin{equation*}
\Phi(z) = \left(\frac{1}{z - \lambda} + \frac{1}{\lambda}\right)^{N} = z^{N}\Phi_{0}(z),
\end{equation*}
则 $\Phi_{0}(z)$ 在 $\widehat{T}(\mathbf{Z})$ 的邻域上有界,故作用于 $M_{22}$.因此,令 $\widehat{S}(n)=\overline{\Phi(\widehat{T}(n))}$($n\in\mathbf{Z}$),则 $S\in M_{p2}\subset M_{pq}$.又由 $\widehat{S}(n)(\widehat{T}(n)-\lambda)^{N}\lambda^{N}=\widehat{T}^{N}(n)$($n\in\mathbf{Z}$),$S(T-\lambda)^{N}\lambda^{N}=T^{N}$.从而,取满足 $h(T)=\lambda$ 的 $h\in\Delta(\mathfrak{M}_{pq})$,与 $\lambda\neq0$ 矛盾.设 $0\notin\overline{\widehat{T}(\mathbf{Z})}$.此时,存在 $\varepsilon_{0}>0$ 使得 $\left|\widehat{T}(n)\right|\geq\varepsilon_{0}>0$($n\in\mathbf{Z}$).因此 $\left|\widehat{T^{N}}(n)\right|\geq\varepsilon_{0}^{N}$($n\in\mathbf{Z}$),对任意三角多项式 $f$,
\begin{equation*}
\varepsilon_{0}^{N}\|f\|_{2} \leq \|T^{N}f\|_{2} \leq C\|f\|_{p}.
\end{equation*}
这与 $L^{2}(\mathbf{T})\neq L^{p}(\mathbf{T})$ 矛盾.
\end{proof}

\begin{note}\label{note:3.5.27}
(1) 设 $2 \leq p < q < \infty$.关于 $M(p, q)$ 的作用函数,可证明与定理~\ref{thm:3.5.25} (1) 类似的事实.

(2) 当 $1 < p < q < 2$ 时,定理~\ref{thm:3.5.25} (1) 的条件是否也是必要条件,是未解决的问题.

(3) 设 $1 < q < p < \infty$.由 Doss \cite{ch3:16} 的结果,与 3.5.4 小节的结果合并可知,$M(p, q)$ 的在 $[-1, 1]$ 上定义的作用函数与整函数一致.
\end{note}

其次转入定理~\ref{thm:3.5.25} 的证明.利用若干引理,分若干阶段证明.下面是省略证明,作为 Stein 插值定理广为人知的事实(cf. \cite{ch3:4} 及其参考文献).

\begin{lemma}[Stein]\label{lem:3.5.28}
设 $1 \leq p_{0}, q_{0}, p_{1}, q_{1} \leq \infty$,$\{T_{z}\}$ 为满足下列性质的算子族.

(1) $z \mapsto \int_{-\pi}^{\pi}T_z(f)g\dd x$($f, g \in L^1(\mathbf{T})$)在 $0 < \operatorname{Re} z < 1$ 上解析,在 $0 \leq \operatorname{Re} z \leq 1$ 上连续.

(2) 存在某个 $0 < a < \pi$,使得
\begin{equation*}
e^{-a|\operatorname{Im}z|}\log\left|\int_{-\pi}^{\pi}T_{z}(f)g\dd x\right|
\end{equation*}
在 $0 \leq \operatorname{Re}z \leq 1$ 上一致有界.

(3) 存在常数 $M_0, M_1$,对任意实数 $y$,$\|T_{iy}f\|_{q_0} \leq M_0\|f\|_{p_0}$($f \in L^{p_0}(\mathbf{T})$),$\|T_{1+iy}f\|_{q_1} \leq M_1\|f\|_{p_1}$($f \in L^{p_1}(\mathbf{T})$)($y$ 是实数)成立.

此时,设 $0 < \theta < 1$,$\operatorname{Re}z = \theta$,$1/p = (1 - \theta)/p_0 + \theta/p_1$,$1/q = (1 - \theta)/q_0 + \theta/q_1$,则 $\|T_\theta f\|_q \leq M_0^{1-\theta}M_1^\theta\|f\|_p$($f \in L^p(\mathbf{T})$)成立.
\end{lemma}

\begin{lemma}\label{lem:3.5.29}
设 $1 < p < q < 2$.设 $T \in \mathfrak{M}_{pq}$ 且 $\widehat{T}(n)$ 取实数值.进一步设 $1/q_{n} = 1/p - n(1/p - 1/q)$,$1/s_{n} = 1/2 + n(1/p - 1/q)$.此时,取定理~\ref{thm:3.5.25}(1) 的 $n_{0}$,对 $n = 1, 2, \ldots, n_{0}$,$T^{n} \in \mathfrak{M}_{s_{n}2}$,且 $\|T^{n}\|_{\mathfrak{M}_{s_{n}2}} \leq \|T\|_{\mathfrak{M}_{pq}}^{n}$ 成立.
\end{lemma}

\begin{proof}
设 $T \in \mathfrak{M}_{pq}$,$\widehat{T}(n)$ 取实数值.此时 $\|T\|_{\mathfrak{M}_{pq}} = \|T\|_{\mathfrak{M}_{q'p'}}$.利用这一点与 Riesz--Thorin 插值定理,令 $1/u = (1-\theta)/p + \theta/q'$,$1/v = (1-\theta)/q + \theta/p'$($0 \leq \theta \leq 1$),则 $T \in M_{uv}$ 且 $\|T\|_{M_{uv}} \leq \|T\|_{M_{pq}}$ 成立.令 $\theta = n(1/p - 1/q)/(1/p + 1/q - 1)$,则 $T \in M_{q_n q_{n+1}}$.由此,$T^n \in M_{q_0 q_1} M_{q_1 q_2} \cdots M_{q_{n+1} q_n}$(其中 $q_0 = p$),且 $\|T^n\|_{\mathfrak{M}_{pq_n}} = \|T\|_{\mathfrak{M}_{pq}}^n$($n = 1, \ldots, n_0$),所以由 $q_{n_0 + 1} \geq 2 > q_{n_0}$ 以及 $\|T^n\|_{\mathfrak{M}_{pq_n}} = \|T^n\|_{\mathfrak{M}_{q_n'p'}}$,$q_n < 2$,令 $1/s_n = (1 - \tau)/p + \tau/q_n'$($1/2 = (1 - \tau)/q_n + \tau/p'$),得 $T^n \in M_{s_n2}$.
\end{proof}

\begin{lemma}\label{lem:3.5.30}
设 $1 < p \leq 2 \leq q \leq p'$.取 $\{a(n)\} \in \ell^{\infty}$,$T \in M_{pq}$,令 $\widehat{T}a(n) = \widehat{T}(n)a(n)$,则下式成立:

(1) $\|Ta\|_{\mathfrak{M}_{p2}} \leq \|\{a(n)\}\|_{\ell^{\infty}}\|T\|_{\mathfrak{M}_{pq}}$.

(2) $\|T^2a\|_{\mathfrak{M}_{pp'}} \leq \|\{a(n)\}\|_{\ell^{\infty}}\|T\|_{\mathfrak{M}_{pq}}^{2}$.
\end{lemma}

\begin{proof}
(1) 设 $f \in L^{p}(\mathbf{T})$,则
\begin{equation*}
\|Taf\|_{2} \leq \|\{a(\cdot)\}\|_{\ell^{\infty}}\|Tf\|_{2} \leq \|\{a(\cdot)\}\|_{\ell^{\infty}}\|T\|_{\mathfrak{M}_{p2}}\|f\|_{p}.
\end{equation*}
这里,由 $\|T\|_{M_{p2}} \leq \|T\|_{M_{pq}}$ 得 $\|Ta\|_{M_{p2}} \leq \|\{a(n)\}\|_{\ell^{\infty}}\|T\|_{M_{pq}}$.

(2) 取 $f \in L^{p}(\mathbf{T})$,
\begin{equation*}
\|T^{2}af\|_{p'} \leq \|T\|_{\mathfrak{M}_{2p'}}\|Taf\|_{2} \leq \|T\|_{\mathfrak{M}_{2p'}}\|Ta\|_{\mathfrak{M}_{p2}}\|f\|_{p}.
\end{equation*}
这里,利用 $\|T\|_{M_{2p^{\prime}}}=\|T\|_{M_{p2}}$ 与 (1),得 $\|T^{2}af\|_{p'}\leq\|\{a(n)\}\|_{\ell^{\infty}}\|T\|_{M_{pq}}^{2}$.
\end{proof}

\begin{proof}[定理~\ref{thm:3.5.25} 的证明]
(1) 不妨设 $\Phi_0$ 是有界函数.设 $1 < p < q < 2$ 且 $n_0 \geq 2$.此时不妨设 $\alpha_1 = \cdots = \alpha_{n_0} = 0$,故设 $\Phi(t) = |t|^{\beta_0 + 1}\Phi_0(t)$.现在,对 $0 \leq \operatorname{Re} z \leq 1$ 定义算子 $S^z$ 为
\begin{equation*}
\widehat{S^z}(n) = \operatorname{sgn}\widehat{T}(n)\left|\widehat{T}(n)\right|^{(n_0-1)z+1}\Phi_{0}(\widehat{T}(n)).
\end{equation*}
此时,
\begin{equation*}
\|S^{z}\|_{\mathfrak{M}_{s_{1}2}} \leq \|\Phi_{0}\|_{\infty}\|T\|_{\mathfrak{M}_{s_{1}2}} \leq \|\Phi_{0}\|_{\infty}\|T\|_{\mathfrak{M}_{pq}} \quad (\operatorname{Re}z = 0)
\end{equation*}
以及
\begin{equation*}
\|S^{z}\|_{\mathfrak{M}_{s_{n_0}2}} \leq \|\Phi_{0}\|_{\infty}\|T^{n_{0}}\|_{\mathfrak{M}_{s_{n_0}2}} \leq \|\Phi_{0}\|_{\infty}\|T\|_{\mathfrak{M}_{pq}}^{n_{0}} \quad (\operatorname{Re}z = 1)
\end{equation*}
成立.由此利用 Stein 插值定理,令 $\theta$ 满足 $1/q = (1-\theta)/s_{1} + \theta/s_{n_{0}}$,则
\begin{equation*}
\|S^{\theta}\|_{\mathfrak{M}_{q2}} \leq \|\Phi_{0}\|_{\infty}\|T\|_{\mathfrak{M}_{pq}}^{1+\theta(n_0-1)}.
\end{equation*}
注意这里的 $\theta$ 取法满足 $\theta(n_{0}-1)+1=\beta_{0}$.换言之,$\widehat{S^{\theta+1}}(n)=\left|\widehat{T}(n)\right|^{\beta_{0}+1}\Phi_{0}(\widehat{T}(n))$,
\begin{equation*}
\|S^{\theta+1}\|_{\mathfrak{M}_{pq}} \leq \|S^{\theta+1}\|_{\mathfrak{M}_{p2}} \leq \|T\|_{\mathfrak{M}_{pq}}\|S^{\theta}\|_{\mathfrak{M}_{q2}} \leq \|\Phi_{0}\|_{\infty}\|T\|_{\mathfrak{M}_{pq}}^{\beta_{0}+1}.
\end{equation*}
这样就证明了 $n_{0} \geq 2$ 的情形.$n_{0} = 1$ 的情形也同样证明.事实上,令 $\widehat{R^{z}}(n) = \operatorname{sgn}\widehat{T}(n)|\widehat{T}(n)|^{z}\Phi_{0}(\widehat{T}(n))$,由 $\|R^{iy}\|_{\mathfrak{M}_{22}} \leq \|\Phi_{0}\|_{\infty}$,$\|R^{1+iy}\|_{\mathfrak{M}_{s_{1}2}} \leq \|\Phi_{0}\|_{\infty}\|T\|_{\mathfrak{M}_{pq}}$,令 $1/q = (1-\theta)/2 + \theta/s_{1}$,则 $\|R^{\theta}\|_{\mathfrak{M}_{q2}} \leq \|\Phi_{0}\|_{\infty}\|T\|_{\mathfrak{M}_{pq}}^{\theta}$.此时注意 $\theta = \beta_{0}$.其余与 $n_{0} \geq 2$ 的情形同样证明.

(2) 设 $T \in \mathfrak{M}_{pq}$,$\widehat{T}(\mathbf{Z}) \subset [-1, 1]$.因 $\|T\|_{\mathfrak{M}_{pq}} = \|T\|_{\mathfrak{M}_{q'p'}}$(cf. \cite{ch3:9}),只需考虑 $p' \geq q$ 的情形.设 $\Phi(t) = |t|^{\beta_1 + 1}\Phi_0(t)$.定义算子 $R^z$ 为 $\widehat{R^z}(n) = \operatorname{sgn}\widehat{T}(n)|\widehat{T}(n)|^z\Phi_0(\widehat{T}(n))$,把引理~\ref{lem:3.5.30} 用于 $a(n) = \widehat{R^z}(n)$,得
\begin{equation*}
\|R^{z}\|_{\mathfrak{M}_{p2}} \leq \|\Phi_{0}\|_{\infty}\|T\|_{\mathfrak{M}_{pq}} \quad (\operatorname{Re}z = 0)
\end{equation*}
以及
\begin{equation*}
\|R^{z}\|_{\mathfrak{M}_{pp'}} \leq \|\Phi_{0}\|_{\infty}\|T\|_{\mathfrak{M}_{pq}}^{2} \quad (\operatorname{Re}z = 1).
\end{equation*}
利用 Stein 插值定理,令 $1/q = (1 - \theta)/2 + \theta/p'$($\theta = \beta_1$),则得 $\|TR^{\theta}\|_{\mathfrak{M}_{pq}} \leq \|\Phi_0\|_\infty\|T\|_{\mathfrak{M}_{pq}}^{\theta+1}$.这样就证明了 (2).
\end{proof}

其次证明定理~\ref{thm:3.5.25} 的 (2) 是必要条件.此时,先注意:由 $\Phi$ 作用于 $M(p,q)$ 自然可导出 $\Phi(0)=0$.

\begin{theorem}\label{thm:3.5.31}
设 $1 < p < 2 \leq q < \infty$,$\Phi$ 为 $[-1, 1]$ 上的函数.若 $\Phi$ 作用于 $M(p,q)$,则 $\Phi(t) = \alpha t + |t|^{\beta_{1}+1}\Phi_{0}(t)$.其中 $\alpha$ 是复数,$\Phi_{0}$ 是有界函数,$\beta_{1}$ 是定理~\ref{thm:3.5.25} (2) 中的量.
\end{theorem}

\begin{proof}
如引理~\ref{lem:3.5.30} 那样,不妨设 $1/2 - 1/q \leq 1/p - 1/2$.如下分若干阶段证明.

第 1 阶段.若 $\Phi$ 作用于 $M(p,q)$,则存在两个正数 $C$ 和 $\eta$,使得当 $\|T\|_{\mathfrak{M}_{pq}} < \eta$ 时 $\|\Phi \circ T\|_{\mathfrak{M}_{pq}} \leq C$ 成立.

事实上,假设对任意自然数 $m$,存在 $T_{m} \in \mathfrak{M}_{pq}$,$\widehat{T_{m}}(\mathbf{Z}) \subset [-1, 1]$,使得 $\|T_{m}\|_{\mathfrak{M}_{pq}} < 1/m^{2}$,$\|\Phi \circ T_{m}\|_{\mathfrak{M}_{pq}} > m$.此时不妨设 $\supp\widehat{T_{m}}$ 是有限集,即 $T_{m}$ 是三角多项式.设 $T_{m}$ 为 $N_{m}$ 次三角多项式.此时,令 $T(x) = \sum_{m=1}^{\infty}e^{in_{m}x}T_{m}(x)$,取 $\{n_{m}\}_{m=1}^{\infty}$ 满足 $n_{m} + 3N_{m} < n_{m+1} - 3N_{m+1}$.由 $T_{m}$ 的构造,$T \in M_{pq}$,故 $\Phi(T) \in M_{pq}$.于是,与 3.3 节的 $V_{\varepsilon}$ 同样地,把 $V_{N_{m}}$ 定义为以全体实数为定义域的函数.利用这个 $V_{N_{m}}$,令 $H_{m}(x) = e^{in_{m}x}V_{N_{m}}(x)$,则 $((\Phi \circ T)H_{m})(x) = e^{in_{m}x}(\Phi \circ T_{m})(x)$.由此得 $\|\Phi \circ T_{m}\|_{M_{pq}} \leq \|H_{m}\|_{1}\|\Phi \circ T\|_{M_{pq}} \leq 3\|\Phi(T)\|_{M_{pq}}$.从而这与 $\{T_{m}\}$ 的取法矛盾.

这里,若 $\Phi$ 作用于 $M(p,q)$,则奇函数 $\Phi(t)-\Phi(-t)$ 与偶函数 $\Phi(t)+\Phi(-t)$ 也作用.因 $\Phi$ 是这两个函数的和,下面讨论这两种情形.

第 2 阶段.证明 $\Phi$ 为偶函数时满足定理的条件.为此,先设如下引理.

\begin{lemma}\label{lem:3.5.32}
对素数 $r(\geq 2)$,存在取值 $\delta^{(r)}(n)=r-1$ 或 $-1$ 的数列 $\{\delta^{(r)}(n)\}$,使得
\begin{equation*}
\left|\sum_{n=1}^{N}\delta^{(r)}(n)e^{inx}\right| < (r-1)r(1+\sqrt{r})\sqrt{N} \quad (0 \leq x \leq 2\pi,\ N = 1, 2, \ldots).
\end{equation*}
此时令 $\Delta_{N}^{(r)}(x)=\sum_{n=1}^{N}\delta^{(r)}(n)e^{inx}$.$\Delta_{N}^{(2)}(x)$ 被称为 Rudin--Shapiro 多项式.
\end{lemma}

\begin{proof}
该引理分两步证明.

步骤 1(cf. \cite{ch3:41}).设 $r$ 为满足 $r \geq 2$ 的自然数,$\alpha = e^{2\pi i/r}$,适当取 $\varepsilon_{r}(n)$($\varepsilon_{r}(n) = 1, \alpha, \ldots, \alpha^{r-1}$),则
\begin{equation*}
\left|\sum_{n=1}^{N}\varepsilon_{r}(n)e^{in\theta}\right| \leq r(1+\sqrt{r})N^{1/2} \quad (0 \leq \theta \leq 2\pi,\ N \geq 1).
\end{equation*}
事实上,设 $A_{0}, A_{1}, \ldots, A_{r-1}$ 为任意复数,利用 $\alpha^{sr} = 1$($s = 0, 1, \ldots, r-1$),
\begin{equation}\tag{3.5.1}
\sum_{s=0}^{r-1}\left|\sum_{j=0}^{r-1}\alpha^{sj}A_{j}\right|^{2} = r\sum_{j=0}^{r-1}|A_{j}|^{2}
\label{eq:3.5.1}\end{equation}
成立.其次,把多项式 $P_{k}^{s}(z)$($k = 0, 1, \ldots$;$s = 0, 1, \ldots, r-1$)定义为 $P_{0}^{0}(z) = P_{0}^{1}(z) = \cdots = P_{0}^{r-1}(z) = z$,并对 $k = 0, 1, 2, \ldots$ 令 $P_{k+1}^{s}(z) = \sum_{j=0}^{r-1}z^{jr^{k}}\alpha^{sj}P_{k}^{j}(z)$($s = 0, 1, \ldots, r-1$),则 $P_{k}^{s}(z)$ 是 $r^{k}$ 次多项式.此时,把 $\varepsilon_{r}(n)$ 取为 $P_{k}^{0}(e^{ix})$ 的第 $n$ 个 Fourier 系数,则知它具有所求的性质.事实上,如下证明.由 $P_{k}^{s}(z)$ 的构造,当 $z = e^{ix}$ 时,利用 \eqref{eq:3.5.1},
\begin{align*}
\sum_{s=0}^{r-1}\left|P_{k+1}^{s}(z)\right|^{2} &= \sum_{s=0}^{r-1}\left|\sum_{j=0}^{r-1}z^{jr^{k}}\alpha^{sj}P_{k}^{j}(z)\right|^{2}\\
&= r\sum_{j=0}^{r-1}\left|z^{jr^{k}}P_{k}^{j}(z)\right|^{2}\\
&\vdots\\
&= r^{k+2}.
\end{align*}
因此
\begin{equation}\tag{3.5.2}
\left|P_{k}^{s}(e^{ix})\right| \leq r^{1/2}r^{k/2} \quad (k = 0, 1, \ldots;\ s = 0, 1, \ldots, r-1)
\label{eq:3.5.2}\end{equation}
得证.特别地,当 $s=0$ 时,$\sum_{n=1}^{r^{k}}\varepsilon_{r}(n)e^{inx}=P_{k}^{0}(e^{ix})$,若 $N=r^{k}$,则步骤 1 的结果得证.

其次,当 $1 \leq n \leq r^{k}$ 时,令 $S_{n}(P_{k}^{s})(x) = \sum_{|m| \leq n}\widehat{P_{k}^{s}}(m)e^{imx}$,则
\begin{equation}\tag{3.5.3}
\left|S_{n}(P_{k}^{s})\right| \leq \sqrt{r}(1+\sqrt{r})r^{k/2} \quad (s = 0, 1, \ldots, r-1)
\label{eq:3.5.3}\end{equation}
可证明.这是通过对 $k$ 的数学归纳法,利用 \eqref{eq:3.5.2} 证明的.从而对 $N \geq 1$,取满足 $r^{k-1} \leq N \leq r^{k}$ 的 $k$,由 $\left|\sum_{n=1}^{N}\varepsilon_{r}(n)e^{inx}\right| = \left|S_{n}(P_{k}^{0})\right|$,从 \eqref{eq:3.5.3} 得
\begin{equation*}
\left|\sum_{n=1}^{N}\varepsilon_{r}(n)e^{in\theta}\right| \leq r(1+\sqrt{r})N^{1/2}.
\end{equation*}

步骤 2.叙述引理~\ref{lem:3.5.32} 的证明.在步骤 1 中,取 $\alpha^{t}$($t=1,2,\ldots,r-1$)代替 $\alpha$,因 $r$ 是素数,$\alpha^{t}$ 是 1 的 $r$ 次根.把 $\alpha^{t}$ 作为 $\alpha$ 按步骤 1 那样构造多项式 $P_{k}^{0}$,则 $\widehat{P_{k}^{0}}(n)=(\varepsilon_{r}(n))^{t}$.这里,令
\begin{equation*}
\delta^{(r)}(n) = \sum_{t=1}^{r-1}(\varepsilon_{r}(n))^{t},
\end{equation*}
则 $\delta^{(r)}(n)=r-1$ 或 $-1$.从而
\begin{align*}
\left|\sum_{n=1}^{N}\delta^{(r)}(n)e^{inx}\right| &= \left|\sum_{n=1}^{N}\left(\sum_{t=1}^{r-1}(\varepsilon_{r}(n))^{t}\right)e^{inx}\right|\\
&\leq \sum_{t=1}^{r-1}\left|\sum_{n=1}^{N}(\varepsilon_{r}(n))^{t}e^{inx}\right|\\
&\leq (r-1)r(1+\sqrt{r})N^{1/2}.
\end{align*}
\end{proof}

\begin{lemma}\label{lem:3.5.33}
利用引理~\ref{lem:3.5.32} 的记号,
\begin{equation*}
\|\Delta_{N}^{(r)}\|_{\mathfrak{M}_{pq}} \leq (r-1)r(1+\sqrt{r})N^{1/p - 1/2}
\end{equation*}
成立.
\end{lemma}

\begin{proof}
由乘子的性质,$\|\Delta_{N}^{(r)}\|_{\mathfrak{M}_{22}}=\|\{\delta^{(r)}(n)\}\|_{l^{\infty}}\leq r-1$,$\|\Delta_{N}^{(r)}\|_{\mathfrak{M}_{1s}}\leq\|\Delta_{N}^{(r)}\|_{s}\leq\|\Delta_{N}^{(r)}\|_{\infty}\leq(r-1)r(1+\sqrt{r})\sqrt{N}$ 成立,故取 $0<\theta<1$ 使 $1/p=\theta/1+(1-\theta)/2$,再取 $s$ 使 $1/q=\theta/s+(1-\theta)/2$,利用 Riesz--Thorin 插值定理即可.
\end{proof}

第 2 阶段的证明.利用引理~\ref{lem:3.5.33},取 $r = 2$,$\Delta_N^{(2)}N^{-(1/p - 1/2)}$ 在 $\mathfrak{M}_{pq}$ 中(在 $N$ 无关的意义下)一致有界.从而,存在常数 $C > 0$ 使得
\begin{equation*}
\left\|\Phi \circ \left(\frac{\Delta_{N}^{(2)}}{C N^{1/p - 1/2}}\right)\right\|_{\mathfrak{M}_{pq}} \leq C.
\end{equation*}
由此,因 $\Phi$ 是偶函数,
\begin{equation*}
\left\|\sum_{n=1}^{N}\Phi\left(\frac{1}{C N^{1/p - 1/2}}\right)e^{inx}\right\|_{q} \leq C\left\|\sum_{n=1}^{N}e^{inx}\right\|_{q}.
\end{equation*}
因此
\begin{equation*}
\left|\Phi\left(\frac{1}{C N^{1/p - 1/2}}\right)\right|N^{1 - 1/q} \leq C N^{1 - 1/p}.
\end{equation*}
此时,把 $\Phi(t)$ 换成 $\Phi(\eta t)$($1/2 < \eta < 1$)也同理可行,所以在 0 的邻域上存在常数 $C > 0$ 使得
\begin{equation*}
|\Phi(t)t^{-1}| \leq C|t|^{\beta_{1}}.
\end{equation*}
这里,令 $\Phi_{0}(t)=\Phi(t)t^{-(1+\beta_{1})}$($t\neq0$),由 $p>1$,$\Phi_{0}$ 有界.事实上,若 $t_{n}\in[-1,1]$,$|\Phi_{0}(t_{n})|\to\infty$,取 $m(k)=t_{n}$($k=2^{n}$),$m(k)=0$($k\neq2^{n}$)的 $m=\{m(k)\}$,由 $\Phi(m(\cdot))\in M(p,q)\subset M(2,2)$,$\Phi(m(\cdot))\in\ell^{\infty}$,故得矛盾(cf. \cite{ch3:11}, \cite{ch3:26}).又,若 $\Phi(t)$ 是奇函数,用同样的方法可证明存在常数 $C>0$ 使得 $|\Phi(t)|\leq C|t|$.这个结果后面使用.以上证明了偶函数时 $\Phi$ 作用于 $M(p,q)$ 的必要条件.

第 3 阶段.若 $\Phi$ 是奇函数,则
\begin{equation*}
\left|\Phi(t) - 2\Phi\left(\frac{t}{2}\right)\right| \leq C|t|^{1 + \beta_{1}}.
\end{equation*}
事实上,如下证明.令
\begin{align*}
\Phi\left(\frac{t}{2}\delta^{(3)}(n)\right) &= \frac{1}{3}\left[\Phi(t) - 2\Phi\left(\frac{t}{2}\right)\right] + \frac{1}{3}\left[\Phi(t) + \Phi\left(\frac{t}{2}\right)\right]\delta^{(3)}(n)\\
&= \Psi_{1}(t) + \Psi_{2}(t)\delta^{(3)}(n).
\end{align*}
此时,由第 1 阶段,引理~\ref{lem:3.5.33} 与第 2 阶段后半的注意,不妨设
\begin{equation*}
\left\|\Phi \circ \left(\frac{\Delta_{N}^{(3)}}{C N^{1/p - 1/2}}\right)\right\|_{\mathfrak{M}_{pq}} \leq C, \quad |\Phi(t)| \leq C|t|.
\end{equation*}
从而,注意 $\Psi_{2}$ 也是作用函数,利用引理~\ref{lem:3.5.33},
\begin{equation*}
\begin{aligned}
\left\|\sum_{n=1}^{N}\Psi_{2}\left(\frac{1}{C N^{1/p - 1/2}}\right)\delta^{(3)}(n)e^{inx}\right\|_{\mathfrak{M}_{pq}} &= \left|\Psi_{2}\left(\frac{1}{C N^{1/p - 1/2}}\right)\right|\|\Delta_{N}^{(3)}\|_{\mathfrak{M}_{pq}} \leq C.
\end{aligned}
\end{equation*}
因此
\begin{equation*}
\left\|\sum_{n=1}^{N}\Psi_{1}\left(\frac{1}{C N^{1/p - 1/2}}\right)e^{inx}\right\|_{M(p,q)} \leq C \quad (N = 1, 2, \ldots).
\end{equation*}
换言之,
\begin{equation*}
\left\|\sum_{n=1}^{N}\Psi_{1}\left(\frac{1}{C N^{1/p - 1/2}}\right)e^{inx}\right\|_{q} \leq C\left\|\sum_{n=1}^{N}e^{inx}\right\|_{p}.
\end{equation*}
从而
\begin{equation*}
\left|\Psi_{1}\left(\frac{1}{C N^{1/p - 1/2}}\right)\right|N^{1 - 1/q} \leq C N^{1 - 1/p},
\end{equation*}
所以
\begin{equation*}
|\Psi_{1}(t)| \leq C t^{\frac{1/p - 1/q}{1/p - 1/2}}.
\end{equation*}
综上所述,第 3 阶段得证.

这里,注意第 3 阶段表明:对 $t \in [-1, 1]$,$\{2^{n}\Phi\left(\frac{t}{2^{n}}\right)\}_{n}$ 是 Cauchy 列.于是令
\begin{equation*}
\Phi_{1}(t) = \lim_{n \to \infty}2^{n}\Phi\left(\frac{t}{2^{n}}\right), \quad \Phi_{2}(t) = \Phi(t) - \Phi_{1}(t).
\end{equation*}

第 4 阶段.$|\Phi_2(t)|\leq C|t|^{1 + \beta_1}$.

第 4 阶段如下证明.利用第 3 阶段的结果,
\begin{align*}
|\Phi_{2}(t)| &= \left|\lim_{n \to \infty}2^{n}\Phi\left(\frac{t}{2^{n}}\right) - \Phi(t)\right|\\
&\leq \sum_{n=1}^{\infty}\left|2^{n}\Phi\left(\frac{t}{2^{n}}\right) - 2^{n-1}\Phi\left(\frac{t}{2^{n-1}}\right)\right|\\
&\leq \sum_{n=1}^{\infty}C'2^{n-1}\left|\frac{t}{2^{n-1}}\right|^{1+\beta_{1}}\\
&\leq C|t|^{1+\beta_{1}}.
\end{align*}
即可证明.这样第 4 阶段得证.

第 5 阶段.$\Phi_{1}(t)$ 是 $[-1, 1]$ 上的连续函数.

事实上,如下证明.必要时把 $\Phi_{1}(t)$ 换成 $\Phi_{1}(\eta t)$,只需证明 $\Phi_{1}(t)$ 在 $t=1$ 处连续.设 $\{t_{n}\}$ 为满足 $|1-t_{m}|<2^{-m}$ 的数列,令
\begin{equation*}
N_{m} = \left[2^{m/(\frac{1}{p} - \frac{1}{2})}\right], \quad T_{m} = \frac{1}{2^{m}}\sum_{n=1}^{N_{m}}\left\{\delta^{(2)}(n) + \frac{1 - t_{m}}{2}(1 - \delta^{(2)}(n))\right\}e^{inx}.
\end{equation*}
此时
\begin{equation*}
\|T_{m}\|_{\mathfrak{M}_{pq}} \leq 2^{-m}\left\|\sum_{n=1}^{N_{m}}\delta^{(2)}(n)e^{inx}\right\|_{\mathfrak{M}_{pq}} + C2^{-2m}\left\|\sum_{n=1}^{N_{m}}e^{inx}\right\|_{\mathfrak{M}_{pq}}.
\end{equation*}
右边第 1 项由 $N_{m}$ 的取法与引理~\ref{lem:3.5.33} 关于 $m$ 一致有界.右边第 2 项可用 $2^{-2m}2^{m/(\frac{1}{p}-\frac{1}{2})\cdot(\frac{1}{p}-\frac{1}{q})}$ 控制,由 $1/p+1/q\geq1$ 关于 $m$ 一致有界(cf. \cite{ch3:40}).可是,因 $\Phi_{1}$ 是奇函数,
\begin{align*}
&\Phi_{1}\left(\delta^{(2)}(n) + \frac{1 - t_{m}}{2}(1 - \delta^{(2)}(n))\right)\\
&\qquad = \frac{1}{2}\left[\Phi_{1}(1) - \Phi_{1}(t_{m})\right] + \frac{1}{2}\left[\Phi_{1}(1) + \Phi_{1}(t_{m})\right]\delta^{(2)}(n).
\end{align*}
利用 $2\Phi_{1}\left(\frac{t}{2}\right)=\Phi_{1}(t)$,若取 $\|T_{m}\|_{\mathfrak{M}_{pq}}\leq2^{m_{0}}$,则从上式
\begin{align*}
C &\geq \|\Phi_{1}(2^{-m_{0}}T_{m})\|_{\mathfrak{M}_{pq}}\\
&= 2^{-m-m_{0}}\|\Phi_{1}(2^{m}T_{m})\|_{\mathfrak{M}_{pq}}\\
&\geq C|\Phi_{1}(1) - \Phi_{1}(t_{m})|2^{-m}N_{m}^{1/p - 1/q} - C2^{-m}N_{m}^{1/p - 1/2},
\end{align*}
且因 $2^{-m}N_{m}^{1/p-1/q}\to\infty$($m\to\infty$),必须 $\Phi_{1}(t_{m})\to\Phi_{1}(1)$($m\to\infty$).因此 $\Phi_{1}(t)$ 在 $t=1$ 处连续.以上第 5 阶段得证.

第 6 阶段.$\Phi_1(t) = \alpha t$($t\in [-1,1]$)成立.

第 6 阶段的证明如下进行.

设 $0 < t < 1$,$r$ 为素数.由 $\Phi_{1}(t)$ 是奇函数,
\begin{equation}\tag{3.5.4}\begin{aligned}
\Phi_{1}\left(\frac{t\delta^{(r)}(n)}{2^{m}(r-1)}\right) &= \frac{1}{2^{m}r}\left[\Phi_{1}(t) - (r-1)\Phi_{1}\left(\frac{t}{r-1}\right)\right]\\
&\quad + \frac{1}{2^{m}r}\left[\Phi_{1}(t) + \Phi_{1}\left(\frac{t}{r-1}\right)\right]\delta^{(r)}(n)
\end{aligned}\label{eq:3.5.4}\end{equation}
成立.可是,由引理~\ref{lem:3.5.33},
\begin{equation*}
\left\|t2^{-m}(r-1)^{-1}\Delta_{N_{m}}^{(r)}\right\|_{M(p,q)} \leq Ct2^{-m}(r-1)^{-1}N_{m}^{1/p - 1/2},
\end{equation*}
故令 $N_{m} = [2^{m/(1/p-1/2)}]$,把自然数 $m_{0}$ 取得足够大,\eqref{eq:3.5.4} 右边第 2 项的算子范数即可用 $2^{m_{0}}$ 控制,
\begin{equation*}
\left\|\Phi_{1}\left(\frac{t\Delta_{N}^{(r)}}{2^{m+m_{0}}(r-1)}\right)\right\|_{\mathfrak{M}_{pq}} \leq C
\end{equation*}
($C$ 可以依赖于 $r$).从而
\begin{align*}
2^{m_{0}}C &\geq \frac{1}{2^{m}r}\left|\Phi_{1}(t) - (r-1)\Phi_{1}\left(\frac{t}{r-1}\right)\right|N_{m}^{1/p - 1/q}\\
&\quad - \frac{C}{2^{m}r}\left|\Phi_{1}(t) + \Phi_{1}\left(\frac{t}{r-1}\right)\right|N_{m}^{1/p - 1/2}.
\end{align*}
因 $2^{-m}N_{m}^{1/p-1/2}$ 有界且 $2^{-m}N_{m}^{1/p-1/q}\to\infty$($m\to\infty$),必须
\begin{equation*}
\Phi_{1}(t) = (r-1)\Phi_{1}\left(\frac{t}{r-1}\right).
\end{equation*}
反复利用这一点,设 $q, r$ 为素数,则
\begin{equation*}
\Phi_{1}\left(\left(\frac{r-1}{q-1}\right)^{n}\right) = \left(\frac{r-1}{q-1}\right)^{n}\Phi_{1}(1).
\end{equation*}
这里,因 $\left\{(\frac{r-1}{q-1})^{n}: r, q \text{ 是素数, } n \text{ 是自然数}\right\}$ 在 $(0, 1)$ 中稠密且 $\Phi_{1}(t)$ 连续,$\Phi_{1}(t)=t\Phi_{1}(1)$($0\leq t\leq1$).以上第 6 阶段得证.

综上所述,全部证明完毕.
\end{proof}

\begin{note}\label{note:3.5.34}
定理~\ref{thm:3.5.31} 的 $q = 2$ 的情形,由 Hatori \cite{ch3:19} 推广到了无限紧 Abel 群上.另外,$M_{pq}$ 的极大理想空间在 Hatori \cite{ch3:20} 中作了研究.
\end{note}

\subsection{\texorpdfstring{$C_0 M_p(\mathbb{Z})$($1<p<2$)的作用函数}{C0 M_p(Z)(1<p<2)的作用函数}}\label{sec:3.5.6}

这里,关于 $\mathfrak{M}_{p}(\mathbf{T})$($1 < p < 2$)的子空间 $C_{0}\mathfrak{M}_{p}(\mathbf{T})$ 的作用函数,只叙述 M. Zafran 的研究 \cite{ch3:59}(cf. \cite{ch3:58})的结果.详情请参阅论文.

\begin{definition}\label{def:3.5.35}
设 $1 < p < 2$.$T \in C_{0}\mathfrak{M}_{p}(\mathbf{T})$ 对应于 $\{\widehat{T}(n)\}_{n \in \mathbf{Z}} \in \ell^{\infty}$(即 $\widehat{T}f(n) = \widehat{T}(n)\widehat{f}(n)$,$f \in L^{2}(\mathbf{T})$),但因不会产生混乱,把 $C_{0}\mathfrak{M}_{p}(\mathbf{T})$ 等同于 $C_{0}M_{p}(\mathbf{Z}) = \{\widehat{T}: T \in C_{0}\mathfrak{M}_{p}(\mathbf{T})\}$.设 $\Phi$ 为 $[-1, 1]$ 上的函数.当 $\widehat{T} \in C_{0}M_{p}(\mathbf{Z})$,$\widehat{T}(\mathbf{Z}) \subset [-1, 1]$ 时,若 $\Phi(\widehat{T}) \in C_{0}M_{p}(\mathbf{Z})$,则称 $\Phi$ 作用于 $C_{0}\mathfrak{M}_{p}(\mathbf{T})$.
\end{definition}

现在,在 3.5.4 小节中证明了满足 $\operatorname{Sp}(\mu, \mathfrak{M}_p(\mathbf{T})) \supsetneqq \overline{\widehat{\mu}(\mathbf{Z})}$ 的 $\mu \in M(\mathbf{T})$ 的存在,而 M. Zafran 注意到:对此 $\mu$,$\hat{\mu}$ 在无穷远点不趋于零.事实上,下式成立.

\begin{theorem}[{$\;$\cite{ch3:22}}]\label{thm:3.5.36}
设 $1 \leq q < p \leq 2$.若 $T \in C_{0}\mathfrak{M}_{q}(\mathbf{T})$,则
\begin{equation*}
\operatorname{Sp}(T, \mathfrak{M}_{p}) = \hat{T}(\mathbf{Z}) \cup \{0\}
\end{equation*}
成立.
\end{theorem}

这是由 Riesz--Thorin 插值定理证明的.关于 3.5.4 小节所述的 Hörmander 问题``$m_p = C_0\mathfrak{M}_p(\mathbf{T})$ ?'',Figa-Talamanca--Gaudry \cite{ch3:18} 给予了否定解决.他们证明:当 $1 < p < 2$ 时,存在 $T_0 \in C_0\mathfrak{M}_p(\mathbf{T})$ 使得 $\|T_0 - T_f\|_{\mathfrak{M}_p} \geq c(>0)$($f \in L^1(\mathbf{T})$).但对这个 $T_0$,Zafran \cite{ch3:57} 证明了 $T_0^2 \in m_p$ 成立,指出 $T_0$ 距 $m_p$ 并不太远.进而,M. Zafran 指出了对应于 $p = 1$ 的情形的 $C_0\mathfrak{M}_1(\mathbf{T}) = M_0(\mathbf{T}) = \{\mu \in M(\mathbf{T}): \hat{\mu} \to 0 \ (|n| \to \infty)\}$ 的作用函数的 Varopoulos 研究 \cite{ch3:51} 的重要性.Varopoulos \cite{ch3:51} 注意到:由 $L^1(\mathbf{T}) \subset M_0(\mathbf{T})$,$M_0(\mathbf{T})$ 的作用函数 $\Phi$ 作用于 $L^1(\mathbf{T})$,故可在 0 的邻域上级数展开;并利用发现 $M_0(\mathbf{T})$ 的特别测度的方法,证明了当 $\Phi(z) = \sum_{n=0}^{\infty}a_nz^n$($|z| < \varepsilon_0$,$\varepsilon_0 > 0$)时,右边级数的收敛半径可以取得任意大.M. Zafran 把 Varopoulos \cite{ch3:51} 的这一方法用于 $C_0\mathfrak{M}_p(\mathbf{T})$ 的作用函数的研究,证明了下面的结果.

\begin{theorem}\label{thm:3.5.37}
设 $1 < p < 2$.$[-1, 1]$ 上的函数 $\Phi$ 作用于 $C_{0}\mathfrak{M}_{p}(\mathbf{T})$ 的充要条件是:$\Phi$ 可延拓为整函数.
\end{theorem}

由这个结果可知,在 $C_0\mathfrak{M}_p(\mathbf{T})$ 中也会发生 Wiener--Pitt 现象.即,

\begin{corollary}\label{cor:3.5.38}
存在 $T \in C_{0}\mathfrak{M}_{p}(\mathbf{T})$ 使得
\begin{equation*}
\operatorname{Sp}\left(T, C_{0}\mathfrak{M}_{p}(\mathbf{T})\right) \supsetneqq \overline{\widehat{T}(\mathbf{Z})}
\end{equation*}
成立.
\end{corollary}

关于定理~\ref{thm:3.5.37} 的证明,经 Brown \cite{ch3:13} 的指摘有所改进,但为得到好的估计,需要使用测度的构造,向量值函数的插值定理,Littlewood--Paley 理论等,比较复杂.最后,关于 $\|e^{i\mu}\|_{\mathfrak{M}_{p}(\mathbf{T})}$ 的估计,叙述 M. Zafran 的结果 \cite{ch3:60}.

\begin{theorem}[Zafran]\label{thm:3.5.39}
设 $1 < p < 2$.此时,对 $j \in Z$,
\begin{equation*}
\sup\left\{\|e^{ij\mu}\|_{\mathfrak{M}_{p}(\mathbf{T})}: \mu \in M(\mathbf{T}),\ \hat{\mu} \in \mathbf{R},\ \|\mu\|_{\mathfrak{M}_{p}(\mathbf{T})} \leq 1\right\} \geq c\beta^{|j|}
\end{equation*}
成立.其中,$c>0$,$\beta>1$ 是不依赖于 $p$ 的常数.
\end{theorem}

与定理~\ref{thm:3.5.21} 或注意~\ref{note:3.5.22} 比较可知,这是体现出 $\mathfrak{M}_{p}(\mathbf{T})$ 与 $M(\mathbf{T})$ 的本质大小差异的估计,是非常有趣的结果.

\section*{结语}

如前言所述,第\ref{ch:3}章主要叙述了据称以 N. Wiener 的定理为开端的谱综合与作用函数的话题.关于谱综合,以 P. Malliavin,N. Th. Varopoulos,S. Saeki,T. W. Körner,D. Müller 等人的研究为中心;关于作用函数,则以 J. P. Kahane,W. Rudin,Y. Katznelson,S. Igari,M. Zafran 等人的研究为中心.本章参考了他们的文献.除这里所选取的内容之外,在与交换 Banach 代数的调和分析相关方面,尚有 \cite{ch3:6},\cite{ch3:15},\cite{ch3:21},\cite{ch3:25},\cite{ch3:30},\cite{ch3:36},\cite{ch3:37},\cite{ch3:38},\cite{ch3:46},\cite{ch3:58} 等大量有趣的研究.另外,参考文献的排列顺序为:先书,后论文.

此番,谨向给予我执笔机会的东京女子大学宮地晶彦教授致以衷心的感谢.
